\section{楽器デバイス（Processing）のプログラムソース}
本節では，楽器デバイスにおける音色表現を担うProcessingプログラムの全文を示す．

\subsection{SoundProcessingCode.pde}
\begin{lstlisting}[caption={SoundProcessingCode.pde}, label={lst:proc_main}]
import ddf.minim.*;
import ddf.minim.analysis.*;
import ddf.minim.effects.*;
//import ddf.minim.signals.*;
//import ddf.minim.spi.*;
import ddf.minim.ugens.*;

//import ddf.minim.*;
//import ddf.minim.ugens.*;
//import ddf.minim.analysis.*;

//import ddf.minim.effects.*; // 追加

import processing.serial.*;
Serial myPort;

Minim minim;
AudioOutput out;
FFT fft;

int musicBPM = 120;

String timbre = "Flute";

byte soundOctave = 60;
byte soundDoReMi = 9;

byte soundStartVelocity = 96;
byte soundEndVelocity = 96;





float currentMasterVolume = 1.0;

// 0:なし　1:フルート　2:ストリングス　3:ピアノ　4:ドラム
final byte arduinoNumber = 3;

boolean isDebugMode = true;



// 反響音用
Summer longOut;
Summer shortOut;


InstrumentPlayer player;
DataUtils utils;








void setup()
{
  size(1024, 850);
  minim = new Minim(this);
  out = minim.getLineOut(Minim.MONO, 1024);
  
  // --- 1. エフェクト専用のバス（部屋）を3つだけ作る ---
  // 1. ストリングス専用のバス（通り道）を生成 (UGen)
  /*longOut = new Summer();
  // 2. 新しいUGen形式の反響音エフェクトを生成 (UGen)
  SimpleConvolutionUGen longReverb = new SimpleConvolutionUGen(20000, 0.17f);
  // 3. 【超重要】すべてを .patch() で数珠繋ぎ（シグナルチェーン）にする
  longOut.patch(longReverb);  // ミキサーの音を、反響音エフェクトへ流す
  longReverb.patch(out); */        // 反響音エフェクトの音を、全体の出口へ流す
  
  longOut = new Summer();
  AddReverb longReverb = new AddReverb(15000, 0.20f);
  longOut.patch(longReverb);
  longReverb.patch(out);
  
  shortOut = new Summer();
  AddReverb shortReverb = new AddReverb(2000, 0.25f);
  shortOut.patch(shortReverb);
  shortReverb.patch(out);
  
  
  player = new InstrumentPlayer(out); // thisとoutを渡す
  
  // ★ここに追記：過去2000サンプル（約45ミリ秒分）を参照する反響音を適用
  //out.patch(new SimpleConvolutionEffect(15000));
  
  //out.setTempo (musicBPM);
  //currentWaveform = Waves.SINE;
  
  utils = new DataUtils();
  // 1. float型の二次元配列を取得
  //float[][] Note_f = utils.GetLUT("Note");
  //// 2. Note_fと同じ大きさで、int型の二次元配列「Note」を作成
  //Note = new int[Note_f.length][Note_f[0].length];
  //// 3. 2重ループで1要素ずつint型に変換して代入
  //for (int i = 0; i < Note_f.length; i++) {
  //  for (int j = 0; j < Note_f[i].length; j++) {
  //    Note[i][j] = int(Note_f[i][j]); // 小数点以下は切り捨て
  //  }
  //}
  
  out.setGain(-6.0f);
  
  fft = new FFT(out.bufferSize(), out.sampleRate());
  
  
  
  audioThread = new Thread(new Runnable() {
    public void run() {
      while (true) {
        if (isPlaying && isDebugMode) {
          long getTime = System.nanoTime();
          while (getTime - lastTime >= interval) {
            SoundPlay();
            lastTime += interval;
            tick++;
          }
        }
        
        
        // CPUの暴走（負荷100%）を防ぐため、1ミリ秒だけ休憩させる
        // これにより、毎秒1000回の頻度で正確にタイマーがチェックされます
        //try {
        //  Thread.sleep(interval / 1000000000 *2); 
        //  //Thread.sleep(1);
        //} catch (InterruptedException e) {
        //  break; // エラーが起きたらループを抜ける
        //}
        
        
        // ➔ 休憩（sleep）を廃止し、CPUに「超高速空回り中」と伝えるだけにする
        //Thread.onSpinWait(); 
        
        
        
        else {
          // 停止中はCPUを休ませる
          try { Thread.sleep(100); } catch (InterruptedException e) {}
        }

        // ➔ sleep(1) を廃止し、これに置き換える
        if (isPlaying) {
          Thread.onSpinWait(); // CPUに「今は待機中だよ」と伝えて効率化する（Java 9以降）
        }
      }
    }
  });
  audioThread.start(); // スレッドをスタート！
  //fpsCheckStartTime = System.nanoTime();
  
  
  // シリアル通信開始処理
  println("=== 利用可能なシリアルポート ===");
  String portList[] = Serial.list();
  println(portList); // ポート一覧を確認
  myPort = new Serial(this, portList[2], 115200); // ポートは適切に変更
  myPort.bufferUntil('\n'); // 開業まで溜めてから
  println("受信待機中・・・");
  println("==========");
  
}



void draw()
{
  background(0);
  //if (!(isPlaying && isDebugMode)) {
  // --- オシロスコープの描画 ---
  stroke(255);
  for (int i=0; i < out.bufferSize() - 1; i++) {
    float x1 = map(i, 0, out.bufferSize(), 0, width);
    float x2 = map(i+1, 0, out.bufferSize(), 0, width);
    line(x1, 150 - out.left.get(i)*250, x2, 150 - out.left.get(i+1)*250);
  }
  
  // --- スペクトルアナライザの描画 ---
  fft.forward(out.mix);
  stroke(255, 0, 255);
  
  for(int i = 0; i < fft.specSize() - 1; i++) {
    // FFTのインデックスが何Hzに相当するか取得
    float f1 = fft.indexToFreq(i);
    float f2 = fft.indexToFreq(i+1);
    
    // 表示範囲(20〜20000Hz)から外れる場合は線を描画しない
    if (f2 < 20 || f1 > 20000) continue;
    
    float amp1 = fft.getBand(i);
    float amp2 = fft.getBand(i+1);
    
    float y1 = 800 - amp1 * 4;
    float y2 = 800 - amp2 * 4;
    
    // 専用関数を使ってX座標を計算
    float x1 = FreqToX(f1);
    float x2 = FreqToX(f2);
    
    line(x1, y1, x2, y2);
  }
  
  // --- 周波数の目盛りとラベルの描画 (X軸) ---
  fill(255);
  textSize(12);
  textAlign(CENTER, TOP);
  
  // 添付画像通りの目盛りとラベルの配列
  float[] labelFreqsX = {20, 30, 40, 50, 60, 80, 100, 200, 300, 400, 500,
                         800, 1000, 2000, 3000, 4000, 6000, 8000, 10000, 20000};
  String[] labelStrsX = {"20", "30", "40", "50", "60", "80", "100", "200", "300", "400", "500",
                         "800", "1k", "2k", "3k", "4k", "6k", "8k", "10k", "20k"};
  
  for (int i = 0; i < labelFreqsX.length; i++) {
    float x = FreqToX(labelFreqsX[i]);
    
    stroke(100); // 縦線（薄いグレー）
    line(x, 300, x, 820);
    
    fill(200);
    text(labelStrsX[i], x, 830); 
  }
  
  // --- 振幅の目盛りとラベルの描画 (Y軸) ---
  textAlign(RIGHT, CENTER); // テキストを右揃えにして数値を綺麗に並べる
  
  // 振幅(amp)を0から120まで、20刻みで横線を描画します (上限を160から120に変更)
  for (int amp = 20; amp <= 120; amp += 20) {
    // スペクトルアナライザの描画と同じ計算式を使用
    float y = 800 - amp * 4;
    
    // オシロスコープの描画領域（Y=300より上）に被らない場合のみ描画
    if (y >= 320) {
      stroke(50); // 横線（暗いグレー）
      line(40, y, width, y); // 数値ラベルと被らないようにX=40から開始
      
      fill(200);
      text(amp, 35, y); // 左端に振幅の数値を描画
    }
  }
  
  // 操作説明
  //fill(200);
  //textAlign(LEFT, TOP);
  //text("Press 'p' to play song. Press '1'-'6' to change waveform.", 20, 20);
  //}
  
  
  /*
  if (isPlaying && isDebugMode) {
    long getTime = System.nanoTime();
    // 現在の時間と、前回記録した時間の差分を計算
    // if ではなく while にすることで、遅れた分（数tick分）をこの1フレーム内で一気に処理します
    while (getTime - lastTime >= interval) {
      tick++;             
      SoundPlay();
      lastTime += interval; // すっきりした累積補正の式
    }
  }
  */
  /*
  long a = System.nanoTime();
  fps ++;
  if (fpsCheckStartTime + 1000000000L <= a) {
    fpsCheckStartTime += 1000000000L;
    println(fps);
    fps = 0;
  }
  */
}
//long fpsCheckStartTime;
//int fps = 0;

// --- 周波数を対数スケールのX座標に変換する関数 ---
float FreqToX(float f) {
  float minFreq = 20.0f;
  float maxFreq = 20000.0f;
  
  // 範囲外の値を制限する
  if (f < minFreq) f = minFreq;
  if (f > maxFreq) f = maxFreq;
  
  // log() を使って対数スケールにマッピングする
  return map(log(f), log(minFreq), log(maxFreq), 0, width);
};



//音符数と遅延の関係を調査したときに使用
//int noteChecker = 0;

void serialEvent(Serial p) {
  //println("受信");
  // 改行まで文字列を読み込む
  String inString = p.readStringUntil('\n');
  
  // 何も読み込めなかった場合は処理停止
  if (inString == null) return;
  
  // 末尾の改行コードや余分な空白を削除
  inString = trim(inString); 
  
  // カンマで分割して配列にする
  String[] list = split(inString, ',');
  
  // データが空の場合は処理停止
  if (list.length <= 1) return;
  
  // ヘッダとデータ数
  String header = list[0];
  int listLength = list.length - 1;
  
  //　ヘッダによる分岐
  switch (header) {
    
    // 発音指示
    case "N" :
      if (listLength % 5 != 0) break;
      for (int i = 5; i <= listLength; i += 5) {
        
        float pitch, duration, startVel, endVel, type;
        try {
          pitch    = float(list[i-4]);
          duration = float(list[i-3]) * (2.5f / musicBPM); // 60000÷(BPM*24)
          startVel = float(list[i-2]);
          endVel   = float(list[i-1]);
          type     = int(list[i]);
          println("N受信(" + list[i-4] + ", " + list[i-3] + ", " + list[i-2] + ", " + list[i-1] + ", " + list[i] + ")"); 
        } catch (NumberFormatException e) {
          continue;
        }
        // 数値の場合のみ処理続行
        
        /*if (pitch  == 72) {
          if (noteChecker == 0) noteChecker -= millis();
          else {noteChecker += millis(); println("所要時間: " + noteChecker); noteChecker = 0;}
        }*/
        
        switch (arduinoNumber) {
          case 1 :
            player.PlayFlute(pitch, duration, startVel, endVel, currentMasterVolume);
            break;
          case 2 :
            player.PlayStrings(pitch, duration, startVel, endVel, currentMasterVolume);
            break;
          case 3 :
            if      (type == 0) player.PlayPiano(pitch, duration, startVel, currentMasterVolume);
            else if (type == 1) player.PlayTrumpet(pitch, duration, startVel, endVel, currentMasterVolume);
            else if (type == 2) player.PlayHorn(pitch, duration, startVel, endVel, currentMasterVolume);
            break;
          case 4 :
            int a = int(list[i-4]) % 12;
            if      (a == 0) player.PlayKick(currentMasterVolume, startVel);
            else if (a == 3) player.PlaySnare(currentMasterVolume, startVel, endVel);
            else if (a == 4) player.PlayHiHat(currentMasterVolume, startVel);
            break;
          default: break;
        }
      }
      break;
      
    // BPM変更指示（60~180）
    case "B" :
      if (listLength != 1) break;
      int getBPM;
      try {
        getBPM = int(list[1]);
        println("B受信 (" + list[1] + ")");
      } catch (NumberFormatException e) {
        break;
      }
      // 数値の場合のみ処理続行
      int minimumBPM = 60;
      int maximumBPM = 180;
      musicBPM = constrain(getBPM, minimumBPM, maximumBPM);
      break;
      
    // マスターボリューム変更指示（0~255）
    case "V" :
      if (listLength != 1) break;
      float getVolume;
      try {
        getVolume = float(list[1]);
        println("V受信 (" + list[1] + ")");
      } catch (NumberFormatException e) {
        break;
      }
      // 数値の場合のみ処理続行
      float minimumVolume = 10.0f;
      float maximumVolume = 50.0f;
      currentMasterVolume = map( constrain( getVolume, minimumVolume, maximumVolume), minimumVolume, maximumVolume, 0.2, 1.0);
      break;
      
    default: break;
  }
}








void Play()
{
  switch(timbre){
    case "Flute":
      player.PlayFlute(soundOctave + soundDoReMi, 2.0f, soundStartVelocity, soundEndVelocity, 1.0f);
      break;
    //case "Flure2":
    //  player.PlayFlute(soundOctave + soundDoReMi + 128, 4.0f, soundStartVelocity, soundEndVelocity, 1.0f);
    //  break;
    case "Strings":
      player.PlayStrings(soundOctave + soundDoReMi, 1.0f, soundStartVelocity, soundEndVelocity, 1.0f);
      break;
    case "Drum":
      if (soundDoReMi == 0) {
        player.PlayKick(1.0f, soundStartVelocity);
      } else if (soundDoReMi == 1) {
        player.PlaySnareOld(1.0f, soundStartVelocity);
      } else if (soundDoReMi == 2) {
        player.PlayHiHatOld(1.0f, soundStartVelocity);
      } else if (soundDoReMi == 3) {
        player.PlaySnare(1.0f, soundStartVelocity, soundEndVelocity);
      } else if (soundDoReMi == 4) {
        player.PlayHiHat(1.0f, soundStartVelocity);
      } //else if (soundDoReMi == 5) {
      //  player.PlayCymbal(1.0f, soundStartVelocity);
      //}
      break;
    case "Piano":
      //int a = millis();
      player.PlayPiano(soundOctave + soundDoReMi, 10.0f, soundStartVelocity, 1.0f);
      //a = millis()-a;
      //println("経過時間: " + a);
      break;
    case "Horn":
      player.PlayHorn(soundOctave + soundDoReMi, 1.0f, soundStartVelocity, soundEndVelocity, 1.0f);
      break;
    case "Trumpet":
      player.PlayTrumpet(soundOctave + soundDoReMi, 1.0f, soundStartVelocity, soundEndVelocity, 1.0f);
      break;
  }
}



void keyPressed() {
  if (isDebugMode == false) return;
  switch (key)
  {
  case '1':
    soundOctave = 24;
    break;
  case '2':
    soundOctave = 36;
    break;
  case '3':
    soundOctave = 48;
    break;
  case '4':
    soundOctave = 60;
    break;
  case '5':
    soundOctave = 72;
    break;
  case '6':
    soundOctave = 84;
    break;
  case '7':
    soundOctave = 96;
    break;
    
  case '0':
    soundStartVelocity = 127;
    break;
  case '-':
    soundStartVelocity = 112;
    break;
  case 'o':
    soundStartVelocity = 96;
    break;
  case 'p':
    soundStartVelocity = 80;
    break;
  case 'l':
    soundStartVelocity = 64;
    break;
  case ';':
    soundStartVelocity = 48;
    break;
  case ',':
    soundStartVelocity = 32;
    break;
  case '.':
    soundStartVelocity = 16;
    break;
  case '^':
    soundEndVelocity = 127;
    break;
  case '¥':
    soundEndVelocity = 112;
    break;
  case '@':
    soundEndVelocity = 96;
    break;
  case '[':
    soundEndVelocity = 80;
    break;
  case ':':
    soundEndVelocity = 64;
    break;
  case ']':
    soundEndVelocity = 48;
    break;
  case '/':
    soundEndVelocity = 32;
    break;
  case '_':
    soundEndVelocity = 16;
    break;
    
  case 'a':
    soundDoReMi = 0;
    Play();
    break;
  case 'w':
    soundDoReMi = 1;
    Play();
    break;
  case 's':
    soundDoReMi = 2;
    Play();
    break;
  case 'e':
    soundDoReMi = 3;
    Play();
    break;
  case 'd':
    soundDoReMi = 4;
    Play();
    break;
  case 'f':
    soundDoReMi = 5;
    Play();
    break;
  case 't':
    soundDoReMi = 6;
    Play();
    break;
  case 'g':
    soundDoReMi = 7;
    Play();
    break;
  case 'y':
    soundDoReMi = 8;
    Play();
    break;
  case 'h':
    soundDoReMi = 9;
    Play();
    break;
  case 'u':
    soundDoReMi = 10;
    Play();
    break;
  case 'j':
    soundDoReMi = 11;
    Play();
    break;
    
  case 'z':
  case 'Z':
    timbre = "Flute";
    break;
  //case 'x':
  //case 'X':
  //  timbre = "Flure2";
  //  break;
  case 'c':
  case 'C':
    timbre = "Strings";
    break;
  case 'v':
  case 'V':
    timbre = "Drum";
    break;
  case 'b':
  case 'B':
    timbre = "Piano";
    break;
  case 'n':
  case 'N':
    timbre = "Horn";
    //Oscil sawOsc = new Oscil(37, 0.1, Waves.SQUARE);
    //sawOsc.patch(new HighPassSP(500f, out.sampleRate())).patch(out);
    //sawOsc = new Oscil(199, 0.01, Waves.SQUARE);
    //sawOsc.patch(new HighPassSP(1000f, out.sampleRate())).patch(out);
    //sawOsc = new Oscil(347, 0.01, Waves.SAW);
    //sawOsc.patch(new HighPassSP(1800f, out.sampleRate())).patch(out);
    //sawOsc = new Oscil(433, 0.01, Waves.SQUARE);
    //sawOsc.patch(new HighPassSP(2200f, out.sampleRate())).patch(out);
    //sawOsc = new Oscil(859, 0.01, Waves.SAW);
    //sawOsc.patch(new HighPassSP(5000f, out.sampleRate())).patch(out);
    //sawOsc = new Oscil(2100, 0.01, Waves.SAW);
    //sawOsc.patch(new HighPassSP(9000f, out.sampleRate())).patch(out);
    //sawOsc = new Oscil(71, 0.01, Waves.SQUARE);
    //sawOsc.patch(new HighPassSP(600f, out.sampleRate())).patch(out);
    //sawOsc = new Oscil(47, 0.01, Waves.SAW);
    //sawOsc.patch(new HighPassSP(550f, out.sampleRate())).patch(out);
    //sawOsc = new Oscil(13, 0.01, Waves.SAW);
    //sawOsc.patch(new HighPassSP(300f, out.sampleRate())).patch(out);
    break;
  case 'm':
  case 'M':
    timbre = "Trumpet";
    break;
    
    
    
    
    
    case 'A' : // 酒場でブギウギ
      tick = 1536;
      returnTick = new int[][] {{2880, 1728, 48}, {2784, 2880, 281}, {4512, 1728, 48}};
      returnIndex = 0;
      noteIndex = 29;
      interval = GetInterval( 139 );
      lastTime = System.nanoTime();
      isPlaying = true;
      break;
      
    case 'S' : // かえるのうた（オクターブ５）
      tick = 0;
      returnTick = new int[][] {{1536, 0, 0}};
      returnIndex = 0;
      noteIndex = 0;
      interval = GetInterval( 120*2 );
      lastTime = System.nanoTime();
      isPlaying = true;
      break;
      
    case 'D' : // 王宮のメヌエット（強弱なし）
      tick = 4608;
      returnTick = new int[][] {{5760, 4608, 651}, {5736, 5832, 850}, {6432, 5856, 856}, {6216, 6456, 993}, {6816, 4608, 651}, {5760, 6816, 1046}, {8832, 4608, 651}};
      returnIndex = 0;
      noteIndex = 651;
      interval = GetInterval( 122 );
      lastTime = System.nanoTime();
      isPlaying = true;
      break;
      
    case 'F' : // トルコ行進曲
      tick = 8856;
      returnTick = new int[][] {{9240, 8856, 1399}, {10008, 9240, 1499}, {10392, 10008, 1669}, {10776, 10392, 1783}, {11544, 10776, 1895}, {11544, 10008, 1669}, {10392, 10008, 1669}, {10392, 8856, 1399}, {9240, 8856, 1399}, {10008, 11544, 2111}, {11928, 11544, 2111}, {11904, 11952, 2225}, {13632, 8856, 1399}};
      returnIndex = 0;
      noteIndex = 1399;
      interval = GetInterval( 120 );
      lastTime = System.nanoTime();
      isPlaying = true;
      break;
      
    case 'G' : // 序曲_中盤（製作中用）
      tick = 96;
      returnTick = new int[][] {{1608, 72, 0}};
      returnIndex = 0;
      noteIndex = 0;
      interval = GetInterval( 120 );
      lastTime = System.nanoTime();
      isPlaying = true;
      break;
      
    case 'H' : // 序曲_終盤なし
      tick = 13632;
      //returnTick = new int[][] {{16488, 14952, 2899}}; // ピアノ用
      returnTick = new int[][] {{16488, 14952, 3053}}; // ストリングス用
      //returnTick = new int[][] {{16488, 14952, 2985}}; // フルート用
      returnIndex = 0;
      noteIndex = 2638;
      interval = GetInterval( 120 );
      lastTime = System.nanoTime();
      isPlaying = true;
      break;
      
    case 'J' : // 序曲_ドラム用
      tick = 768;
      returnTick = new int[][] {{3624, 2088, 637}};
      returnIndex = 0;
      noteIndex = 92;
      interval = GetInterval( 120 );
      lastTime = System.nanoTime();
      isPlaying = true;
      break;
      
    case 'Q' : // 演奏停止
      isPlaying = false;
      break;
      
    default:
      break;
  }
}
\end{lstlisting}

\subsection{Utils.pde}
\begin{lstlisting}[caption={Utils.pde}, label={lst:proc_utils}]
class DataUtils {
  private float[][] note;
  
  private float[][][] fluteFFLUT, fluteMFLUT, flutePPLUT;
  
  private float[][][] violinFFLUT, violinMFLUT, violinPPLUT;
  private float[][][] violaFFLUT,  violaMFLUT,  violaPPLUT;
  private float[][][] celloFFLUT,  celloMFLUT,  celloPPLUT;
  private float[][][] bassFFLUT,   bassMFLUT,   bassPPLUT;
  
  private float[][][] pianoFFLUT, pianoMFLUT, pianoPPLUT;
  private float[][] pianoDecayFFLUT, pianoWobbleRateFFLUT, pianoWobbleDepthFFLUT;
  private float[][] pianoDecayMFLUT, pianoWobbleRateMFLUT, pianoWobbleDepthMFLUT;
  private float[][] pianoDecayPPLUT, pianoWobbleRatePPLUT, pianoWobbleDepthPPLUT;
  
  private float[][] kickLUT, snareLUT, hihatLUT;
  
  
  
  private float[][][] hornFFLUT,   hornMFLUT,   hornPPLUT;
  private float[][][] trumpetFFLUT,   trumpetMFLUT,   trumpetPPLUT;
  
  
  DataUtils() {
    //println("Loading JSON data...");
    
    // JSONファイルからデータをロードする
    
    //note = Load2LUT("序曲_フルート冒頭_processing.json");
    //note = Load2LUT("序曲_フルート中盤_processing.json");
    //note = Load2LUT("序曲_ピアノ冒頭_processing.json");
    //note = Load2LUT("序曲_ピアノ中盤_processing.json");
    //note = Load2LUT("序曲_ストリングス冒頭_processing.json");
    //note = Load2LUT("序曲_ストリングス中盤_processing.json");
    //note = Load2LUT("序曲_ドラム冒頭_processing.json");
    //note = Load2LUT("序曲_ドラム中盤_processing.json");
    
    //note = Load2LUT("かえるのうた8小節ver「酒場でブギウギ(強弱付き)」「王宮のメヌエット(強弱なし)」「トルコ行進曲」_processing.json");
    //note = Load2LUT("共通部分0706_processing.json");
    //note = Load2LUT("ドラム0706_1_processing.json");
    //note = Load2LUT("ピアノ0706_1_processing.json");
    note = Load2LUT("ストリングス0706_1_processing.json");
    //note = Load2LUT("フルート0706_1_processing.json");
    
    
    fluteFFLUT = Load3LUT("fluteFFLUT.json");
    fluteMFLUT = Load3LUT("fluteMFLUT.json");
    flutePPLUT = Load3LUT("flutePPLUT.json");
    
    violinFFLUT = Load3LUT("violinFFLUT.json");
    violinMFLUT = Load3LUT("violinMFLUT.json");
    violinPPLUT = Load3LUT("violinPPLUT.json");
    
    violaFFLUT = Load3LUT("violaFFLUT.json");
    violaMFLUT = Load3LUT("violaMFLUT.json");
    violaPPLUT = Load3LUT("violaPPLUT.json");
    
    celloFFLUT = Load3LUT("celloFFLUT.json");
    celloMFLUT = Load3LUT("celloMFLUT.json");
    celloPPLUT = Load3LUT("celloPPLUT.json");
    
    bassFFLUT = Load3LUT("bassFFLUT.json");
    bassMFLUT = Load3LUT("bassMFLUT.json");
    bassPPLUT = Load3LUT("bassPPLUT.json");
    
    pianoFFLUT = Load3LUT("pianoFFLUT.json");
    pianoMFLUT = Load3LUT("pianoMFLUT.json");
    pianoPPLUT = Load3LUT("pianoPPLUT.json");
    pianoDecayFFLUT = Load2LUT("pianoDecaysFF.json");
    pianoDecayMFLUT = Load2LUT("pianoDecaysMF.json");
    pianoDecayPPLUT = Load2LUT("pianoDecaysPP.json");
    pianoWobbleRateFFLUT = Load2LUT("pianoWobbleRatesFF.json");
    pianoWobbleRateMFLUT = Load2LUT("pianoWobbleRatesMF.json");
    pianoWobbleRatePPLUT = Load2LUT("pianoWobbleRatesPP.json");
    pianoWobbleDepthFFLUT = Load2LUT("pianoWobbleDepthsFF.json");
    pianoWobbleDepthMFLUT = Load2LUT("pianoWobbleDepthsMF.json");
    pianoWobbleDepthPPLUT = Load2LUT("pianoWobbleDepthsPP.json");
    
    // ドラムのLUTをロード
    kickLUT  = Load2LUT("kickLUT.json");
    snareLUT = Load2LUT("snareLUT.json");
    hihatLUT = Load2LUT("hihatLUT.json");
    
    
    
    
    
    
    
    
    hornFFLUT = Load3LUT("hornFFLUT.json");
    hornMFLUT = Load3LUT("hornMFLUT.json");
    hornPPLUT = Load3LUT("hornPPLUT.json");
    
    trumpetFFLUT = Load3LUT("trumpetFFLUT.json");
    trumpetMFLUT = Load3LUT("trumpetMFLUT.json");
    trumpetPPLUT = Load3LUT("trumpetPPLUT.json");
    
    //println("Data loaded successfully!");
  }
  
  
  // JSONファイルを読み込んで float[][] 配列に変換する便利関数
  float[][] Load2LUT(String fileName) {
    JSONArray jsonArray = loadJSONArray(fileName);
    float[][] lut = new float[jsonArray.size()][];
    
    for (int i = 0; i < jsonArray.size(); i++) {
      if (jsonArray.isNull(i)) {
        lut[i] = null; // 今までの仕様通り、存在しない音程は null にする
        continue;
      }
      JSONArray row = jsonArray.getJSONArray(i);
      if (row.size() > 0) {
        lut[i] = new float[row.size()];
        for (int j = 0; j < row.size(); j++) {
          lut[i][j] = row.getFloat(j);
        }
      } else {
        lut[i] = null; // データが欠損している部分は null にする（今までの仕様通り）
      }
    }
    return lut;
  }
  
  
  // JSONファイルを読み込んで float[][][] の3次元可変長配列に変換する関数
  float[][][] Load3LUT(String fileName) {
    JSONArray jsonArray = loadJSONArray(fileName);
    // 第1次元（音程の数）を確定
    float[][][] lut = new float[jsonArray.size()][][];
    
    for (int i = 0; i < jsonArray.size(); i++) {
      // JSON上で null になっている、または要素が欠損している場合の処理
      if (jsonArray.isNull(i)) {
        lut[i] = null; // 今までの仕様通り、存在しない音程は null にする
        continue;
      }
      
      JSONArray peaksArray = jsonArray.getJSONArray(i);
      int numPeaks = peaksArray.size();
      
      if (numPeaks > 0) {
        // 第2次元（その音程が持つピークの数）を動的に確定
        lut[i] = new float[numPeaks][2]; 
        
        for (int j = 0; j < numPeaks; j++) {
          JSONArray peakPair = peaksArray.getJSONArray(j);
          
          // [周波数, 振幅] のペア（要素数2）を確実に取得
          if (peakPair.size() == 2) {
            lut[i][j][0] = peakPair.getFloat(0); // 周波数 (Hz)
            lut[i][j][1] = peakPair.getFloat(1); // 相対振幅 (0.0 ~ 1.0)
          }
        }
      } else {
        lut[i] = null; // ピーク数が0個の場合もデータなしとして null を代入
      }
    }
    return lut;
  }
  
  
  
  // ① 汎用版：特定のLUTから、指定したピッチの倍音レシピを安全に取り出す関数
  /*
  float[] GetProfileForPitch(float pitch, float[][] lut, int numHarmonics, int baseMidiNote) {
    
    float[] result = new float[numHarmonics];
    
    // LUTの要素数から自動的に最高音を算出して安全にクランプする
    float maxMidiNote = baseMidiNote + lut.length - 1;
    float clampedPitch = constrain(pitch, baseMidiNote, maxMidiNote);
    
    int baseIndexLower = floor(clampedPitch) - baseMidiNote;
    int baseIndexUpper = ceil(clampedPitch) - baseMidiNote;
    
    // 下に向かって一番近い有効データを探す
    int lowerIndex = baseIndexLower;
    while (lowerIndex >= 0 && lut[lowerIndex] == null) { lowerIndex--; }
    
    // 上に向かって一番近い有効データを探す
    int upperIndex = baseIndexUpper;
    while (upperIndex < lut.length && lut[upperIndex] == null) { upperIndex++; }
    
    
    // ==============================================
    // 限界値のセーフティーネット
    // ==============================================
    // 上に有効データが一つも無かったら、下のデータを採用する（最高音の延長）
    if (upperIndex >= lut.length) upperIndex = lowerIndex;
    
    // 下に有効データが一つも無かったら、上のデータを採用する（最低音の延長）
    if (lowerIndex < 0) lowerIndex = upperIndex; 
    
    if (lowerIndex < 0 && upperIndex >= lut.length) {
      // どちらにもデータが無い場合のフォールバック（基音のみ1.0）
      float[] fallback = new float[numHarmonics];
      fallback[0] = 1.0f;
      return fallback;
    }
    
    float lowerMidi = lowerIndex + baseMidiNote;
    float upperMidi = upperIndex + baseMidiNote;
    
    float[] lowerProfile = lut[lowerIndex];
    float[] upperProfile = lut[upperIndex];
    
    float t = 0;
    if (upperMidi > lowerMidi) {
      t = (clampedPitch - lowerMidi) / (upperMidi - lowerMidi);
    }
    
    for (int i = 0; i < numHarmonics; i++) {
      result[i] = lerp(lowerProfile[i], upperProfile[i], t);
    }
    
    return result;
  }
  */
  
  /*
  // ① 刷新版：特定の3次元LUTから、指定したピッチに「最も近い有効データ」を生のまま取り出す関数
  // 戻り値：float[ピーク数][2] -> [h][0]:周波数(Hz), [h][1]:振幅
  float[][] GetProfileForPitch(float pitch, float[][][] lut, int baseMidiNote) {
    
    // LUTの要素数から自動的に最高音を算出して安全にクランプする
    float maxMidiNote = baseMidiNote + lut.length - 1;
    float clampedPitch = constrain(pitch, baseMidiNote, maxMidiNote);
    
    // 四捨五入(round)で、最もピッチが近いサンプルのインデックスを狙う
    int targetIndex = round(clampedPitch) - baseMidiNote;
    targetIndex = constrain(targetIndex, 0, lut.length - 1);
    
    // もし狙ったインデックスがnull（データ欠損）の場合、最も近い「有効なデータ」を探索する
    int finalIndex = targetIndex;
    if (lut[finalIndex] == null) {
      int lowerIndex = finalIndex;
      while (lowerIndex >= 0 && lut[lowerIndex] == null) { lowerIndex--; }
      
      int upperIndex = finalIndex;
      while (upperIndex < lut.length && lut[upperIndex] == null) { upperIndex++; }
      
      // 上下のうち、存在する方、あるいはより近い方を採用する
      if (lowerIndex >= 0 && upperIndex < lut.length) {
        if ((finalIndex - lowerIndex) <= (upperIndex - finalIndex)) {
          finalIndex = lowerIndex;
        } else {
          finalIndex = upperIndex;
        }
      } else if (lowerIndex >= 0) {
        finalIndex = lowerIndex;
      } else if (upperIndex < lut.length) {
        finalIndex = upperIndex;
      } else {
        // 万が一、LUT全体が空っぽだった場合の完全フォールバック（MIDIピッチから計算した基音のみ生成）
        float f0 = 440.0f * pow(2.0f, (clampedPitch - 69.0f) / 12.0f);
        float[][] fallback = new float[1][2];
        fallback[0][0] = f0;   // 周波数
        fallback[0][1] = 1.0f; // 振幅
        return fallback;
      }
    }
    
    // 見つかった有効なサンプルのデータをディープコピーして返す（参照渡しによるデータ破壊を防ぐため）
    float[][] sourceProfile = lut[finalIndex];
    int numPeaks = sourceProfile.length;
    float[][] result = new float[numPeaks][2];
    
    for (int h = 0; h < numPeaks; h++) {
      result[h][0] = sourceProfile[h][0]; // 周波数 (Hz)
      result[h][1] = sourceProfile[h][1]; // 振幅
    }
    
    return result;
  }
  */
  
  
  // ① 確定修正版：特定の3次元LUTから、指定したピッチに「最も近い有効データ」を取り出し、正確なピッチへシフトして返す関数
  // 戻り値：float[ピーク数][2] -> [h][0]:周波数(Hz), [h][1]:振幅
  float[][] GetProfileForPitch(float pitch, float[][][] lut, int baseMidiNote) {
    
    // 【バグ解決の核心1】ユーザーが要求した本来のピッチ（小数ビブラートや範囲外のC3/A7など）を完全に生のまま保持
    float originalPitch = pitch;
    
    // LUTの要素数から自動的に最高音を算出し、インデックス探索用のみ安全にクランプする
    float maxMidiNote = baseMidiNote + lut.length - 1;
    float searchPitch = constrain(pitch, baseMidiNote, maxMidiNote);
    
    // 四捨五入(round)で、最もピッチが近いサンプルのインデックスを狙う
    int targetIndex = round(searchPitch) - baseMidiNote;
    targetIndex = constrain(targetIndex, 0, lut.length - 1);
    
    // もし狙ったインデックスがnull（データ欠損）の場合、最も近い「有効なデータ」を探索する
    int finalIndex = targetIndex;
    println(finalIndex);
    if (lut[finalIndex] == null) {
      int lowerIndex = finalIndex;
      while (lowerIndex >= 0 && lut[lowerIndex] == null) { lowerIndex--; }
      
      int upperIndex = finalIndex;
      while (upperIndex < lut.length && lut[upperIndex] == null) { upperIndex++; }
      
      // 上下のうち、存在する方、あるいは要求された元のピッチ（searchPitch）により近い方を採用する
      if (lowerIndex >= 0 && upperIndex < lut.length) {
        if (abs(searchPitch - (lowerIndex + baseMidiNote)) <= abs((upperIndex + baseMidiNote) - searchPitch)) {
          finalIndex = lowerIndex;
        } else {
          finalIndex = upperIndex;
        }
      } else if (lowerIndex >= 0) {
        finalIndex = lowerIndex;
      } else if (upperIndex < lut.length) {
        finalIndex = upperIndex;
      } else {
        // 万が一、LUT全体が空っぽだった場合の完全フォールバック（MIDIピッチから計算した基音のみ生成）
        float f0 = 440.0f * pow(2.0f, (originalPitch - 69.0f) / 12.0f);
        float[][] fallback = new float[1][2];
        fallback[0][0] = f0;   // 周波数
        fallback[0][1] = 1.0f; // 振幅
        return fallback;
      }
    }
    println(finalIndex);
    
    // 見つかったサンプルのデータを取得
    float[][] sourceProfile = lut[finalIndex];
    int numPeaks = sourceProfile.length;
    float[][] result = new float[numPeaks][2];
    
    // =================================================================
    // 【バグ解決の核心2：流用元ピッチからの絶対距離によるシフト計算】
    // =================================================================
    // 流用元のサンプルが本来持っている「正しい基準MIDIノート番号」を逆算
    float sourceMidi = finalIndex + baseMidiNote;
    
    // 制限のない本来の要求ピッチ（originalPitch）と、流用元（sourceMidi）のガチの差分（半音単位）を計算。
    // これにより、範囲外のC3を弾いた時（例：48 - 59 = -11半音）でも正確なスケーリング係数が算出されます！
    float pitchShiftFactor = pow(2.0f, (originalPitch - sourceMidi) / 12.0f);
    // =================================================================
    
    for (int h = 0; h < numPeaks; h++) {
      // 元のデータが持っている「実測の生の周波数構造」に対して、計算した倍率を掛け算する
      // これにより、インハーモニシティのデコボコを100%維持したまま、滑らかにピッチがスケーリングされます
      result[h][0] = sourceProfile[h][0] * pitchShiftFactor; 
      result[h][1] = sourceProfile[h][1]; // 振幅はそのままコピー
    }
    
    return result;
  }
  
  
  
  
  
  // ② 汎用版：ピッチとベロシティの2軸から、最終的な倍音レシピを生成する関数
  /*
  float[] GetDynamicProfile(float pitch, float velocity, String type, int numHarmonics, int baseMidiNote, float basePower) {
    float[] profile = new float[numHarmonics];
    
    // 3つのLUTから、現在のピッチに応じた倍音構成をそれぞれ取得
    float[] profileFF, profilePP, profileMF;
    
    switch (type)
    {
    case "Flute":
      profilePP = GetProfileForPitch(pitch, flutePPLUT, numHarmonics, baseMidiNote);
      profileMF = GetProfileForPitch(pitch, fluteMFLUT, numHarmonics, baseMidiNote);
      profileFF = GetProfileForPitch(pitch, fluteFFLUT, numHarmonics, baseMidiNote);
      break;
    case "Violin":
      profilePP = GetProfileForPitch(pitch, violinPPLUT, numHarmonics, baseMidiNote);
      profileMF = GetProfileForPitch(pitch, violinMFLUT, numHarmonics, baseMidiNote);
      profileFF = GetProfileForPitch(pitch, violinFFLUT, numHarmonics, baseMidiNote);
      break;
    case "Viola":
      profilePP = GetProfileForPitch(pitch, violaPPLUT, numHarmonics, baseMidiNote);
      profileMF = GetProfileForPitch(pitch, violaMFLUT, numHarmonics, baseMidiNote);
      profileFF = GetProfileForPitch(pitch, violaFFLUT, numHarmonics, baseMidiNote);
      break;
    case "Cello":
      profilePP = GetProfileForPitch(pitch, celloPPLUT, numHarmonics, baseMidiNote);
      profileMF = GetProfileForPitch(pitch, celloMFLUT, numHarmonics, baseMidiNote);
      profileFF = GetProfileForPitch(pitch, celloFFLUT, numHarmonics, baseMidiNote);
      break;
    case "Bass":
      profilePP = GetProfileForPitch(pitch, bassPPLUT, numHarmonics, baseMidiNote);
      profileMF = GetProfileForPitch(pitch, bassMFLUT, numHarmonics, baseMidiNote);
      profileFF = GetProfileForPitch(pitch, bassFFLUT, numHarmonics, baseMidiNote);
      break;
    case "Piano":
      profilePP = GetProfileForPitch(pitch, pianoPPLUT, numHarmonics, baseMidiNote);
      profileMF = GetProfileForPitch(pitch, pianoMFLUT, numHarmonics, baseMidiNote);
      profileFF = GetProfileForPitch(pitch, pianoFFLUT, numHarmonics, baseMidiNote);
      break;
    default:
      println("Error: LUT type '" + type + "'が見つかりません!! ピアノを適用しました。");
      profilePP = GetProfileForPitch(pitch, pianoPPLUT, numHarmonics, baseMidiNote);
      profileMF = GetProfileForPitch(pitch, pianoMFLUT, numHarmonics, baseMidiNote);
      profileFF = GetProfileForPitch(pitch, pianoFFLUT, numHarmonics, baseMidiNote);
      break;
    }
    
    // ベロシティ(0〜127)を0.0〜1.0に正規化
    float normVel = constrain(velocity, 0, 127) / 127.0f;
    
    // ユーザー定義のベロシティ閾値
    float threshPP = 32.0f / 127.0f;
    float threshMF = 80.0f / 127.0f;
    float threshFF = 112.0f / 127.0f;
    
    for (int i = 0; i < numHarmonics; i++) {
      if (normVel <= threshPP) {
        // 0〜32 (無音〜pp) の間は、波形の形はppのまま固定（音量だけが小さくなる）
        profile[i] = profilePP[i];
        
      } else if (normVel <= threshMF) {
        // 32〜80 (pp〜mf) の間を補間
        float t = map(normVel, threshPP, threshMF, 0.0f, 1.0f);
        profile[i] = lerp(profilePP[i], profileMF[i], t);
        
      } else if (normVel <= threshFF) {
        // 80〜112 (mf〜ff) の間を補間
        float t = map(normVel, threshMF, threshFF, 0.0f, 1.0f);
        profile[i] = lerp(profileMF[i], profileFF[i], t);
        
      } else {
        // 112〜127 (ff〜fff) の間は、波形の形はffのまま固定（音量は最大へ向かう）
        profile[i] = profileFF[i];
      }
    }
    
    // 合計値に対する割合（シェア）による音量算出
    float totalSum = 0;
    for (int i = 0; i < numHarmonics; i++) {
      totalSum += profile[i]; 
    }
    
    // ベロシティに応じた「目標の全体音量（パイの大きさ）」を決定する
    // 音量(0.0)の場合はここでtargetVolumeが0になり、完全に無音になります
    float targetVolume = basePower * pow(normVel, 0.7f);
    
    // 各倍音の割合(シェア)を計算し、目標音量を掛け合わせる
    for (int i = 0; i < numHarmonics; i++) {
      // profile[i] / totalSum が「割合(0.0〜1.0)」。そこにtargetVolumeを掛ける。
      // +0.0001f は、ゼロ除算（0で割ってエラーになること）を防ぐためのおまじないです。
      profile[i] = (profile[i] / (totalSum + 0.0001f)) * targetVolume; 
    }
    
    return profile;
  }
  */
  
  
  // ② 刷新版：ピッチとベロシティの2軸から、最終的な「周波数と音量のペア」の2次元配列を生成する関数
  // 戻り値：float[ピーク数][2] -> [h][0]:周波数(Hz), [h][1]:最終音量
  float[][] GetDynamicProfile(float pitch, float velocity, String type, int baseMidiNote, float basePower) {
    
    // 3つの3次元LUTから、現在のピッチに最も近い実測スペクトルをそれぞれ取得
    float[][] profilePP, profileMF, profileFF;
    
    switch (type) {
      case "Flute":
        profilePP = GetProfileForPitch(pitch, flutePPLUT, baseMidiNote);
        profileMF = GetProfileForPitch(pitch, fluteMFLUT, baseMidiNote);
        profileFF = GetProfileForPitch(pitch, fluteFFLUT, baseMidiNote);
        break;
      case "Violin":
        profilePP = GetProfileForPitch(pitch, violinPPLUT, baseMidiNote);
        profileMF = GetProfileForPitch(pitch, violinMFLUT, baseMidiNote);
        profileFF = GetProfileForPitch(pitch, violinFFLUT, baseMidiNote);
        break;
      case "Viola":
        profilePP = GetProfileForPitch(pitch, violaPPLUT, baseMidiNote);
        profileMF = GetProfileForPitch(pitch, violaMFLUT, baseMidiNote);
        profileFF = GetProfileForPitch(pitch, violaFFLUT, baseMidiNote);
        break;
      case "Cello":
        profilePP = GetProfileForPitch(pitch, celloPPLUT, baseMidiNote);
        profileMF = GetProfileForPitch(pitch, celloMFLUT, baseMidiNote);
        profileFF = GetProfileForPitch(pitch, celloFFLUT, baseMidiNote);
        break;
      case "Bass":
        profilePP = GetProfileForPitch(pitch, bassPPLUT, baseMidiNote);
        profileMF = GetProfileForPitch(pitch, bassMFLUT, baseMidiNote);
        profileFF = GetProfileForPitch(pitch, bassFFLUT, baseMidiNote);
        break;
      case "Piano":
        profilePP = GetProfileForPitch(pitch, pianoPPLUT, baseMidiNote);
        profileMF = GetProfileForPitch(pitch, pianoMFLUT, baseMidiNote);
        profileFF = GetProfileForPitch(pitch, pianoFFLUT, baseMidiNote);
        break;
      case "Horn":
        profilePP = GetProfileForPitch(pitch, hornPPLUT, baseMidiNote);
        profileMF = GetProfileForPitch(pitch, hornMFLUT, baseMidiNote);
        profileFF = GetProfileForPitch(pitch, hornFFLUT, baseMidiNote);
        break;
      case "Trumpet":
        profilePP = GetProfileForPitch(pitch, trumpetPPLUT, baseMidiNote);
        profileMF = GetProfileForPitch(pitch, trumpetMFLUT, baseMidiNote);
        profileFF = GetProfileForPitch(pitch, trumpetFFLUT, baseMidiNote);
        break;
      default:
        println("Error: LUT type '" + type + "'が見つかりません!! ピアノを適用しました。");
        profilePP = GetProfileForPitch(pitch, pianoPPLUT, baseMidiNote);
        profileMF = GetProfileForPitch(pitch, pianoMFLUT, baseMidiNote);
        profileFF = GetProfileForPitch(pitch, pianoFFLUT, baseMidiNote);
        break;
    }
    
    // ベロシティ(0〜127)を0.0〜1.0に正規化
    float normVel = constrain(velocity, 0, 127) / 127.0f;
    
    // ユーザー定義のベロシティ閾値
    float thresholdPP = 32.0f / 127.0f;
    float thresholdMF = 80.0f / 127.0f;
    float thresholdFF = 112.0f / 127.0f;
    
    // 現在のベロシティ層に応じて、ベースとなる波形プロファイル（周波数・振幅の骨組み）を決定
    float[][] baseProfile;
    float t = 0;
    //boolean needInterpolation = false;
    //float[][] lowerProfile = null;
    //float[][] upperProfile = null;
    
    if (normVel <= thresholdPP) {
      baseProfile = profilePP;
    } else if (normVel <= thresholdMF) {
      // pp ～ mf の間：構造が近い方のスペクトル形状を選択（あるいは、より高度に処理するためのフラグ設定）
      t = map(normVel, thresholdPP, thresholdMF, 0.0f, 1.0f);
      baseProfile = (t < 0.5f) ? profilePP : profileMF;
    } else if (normVel <= thresholdFF) {
      // mf ～ ff の間
      t = map(normVel, thresholdMF, thresholdFF, 0.0f, 1.0f);
      baseProfile = (t < 0.5f) ? profileMF : profileFF;
    } else {
      baseProfile = profileFF;
    }
    
    //if (normVel >= thresholdFF) baseProfile = profileFF;
    //else if (normVel >= thresholdMF) baseProfile = profileMF;
    //else baseProfile = profilePP;
    
    // 決定した形状ベース（baseProfile）を元に、最終出力用の2次元配列を確保
    int numPeaks = baseProfile.length;
    int validCount = 0;
    float ampThreshold = 0.05f;
    // 1. まずは有効なデータの数をカウントするだけ（上限は変更しない）
    for (int h = 0; h < numPeaks; h++) { 
      if (baseProfile[h][1] >= ampThreshold) validCount++; 
    }
    // 2. 判明した有効な数で、ぴったりサイズの配列を初期化
    float[][] finalProfile = new float[validCount][2];
    
    
    // 周波数は選択したベースサンプルの値を100%そのまま継承（濁りを防ぐ）
    //float totalSum = 0;
    //for (int h = 0; h < numPeaks; h++) { 
    //  if(baseProfile[h][1] < 0.05f) continue;
    //  finalProfile[h][0] = baseProfile[h][0]; // 周波数 (Hz) 固定
    //  finalProfile[h][1] = baseProfile[h][1]; // 振幅の初期値
    //  //totalSum += finalProfile[h][1];
    //}
    
    
    //// ベロシティに応じた「目標の全体音量（エネルギーの総量）」の計算
    //float targetVolume = basePower * pow(normVel, 0.9f);
    
    //// 全体のエネルギー配分（シェア）を計算し、各ピークに最終音量を割り当てる
    ///*for (int h = 0; h < numPeaks; h++) {
    //  //finalProfile[h][1] = (finalProfile[h][1] / (totalSum + 0.0001f)) * targetVolume;
    //  finalProfile[h][1] = finalProfile[h][1] * targetVolume;
    //}*/
    
    //float energy = 0;
    //for(int h = 0; h < numPeaks; h++){
    //  energy += baseProfile[h][1] * baseProfile[h][1];
    //}
    //float gain = targetVolume / sqrt(energy + 1e-9f);
    
    //for (int h = 0; h < numPeaks; h++) {
    //  /*float freq = baseProfile[h][0]; // 周波数 (Hz)
      
    //  // --- 低音ブーストの計算 ---
    //  // 例：200Hz以下を緩やかに持ち上げるカーブ（ローシェルフ風）
    //  float bassBoost = 1.0f;
    //  if (freq < 400.0f) {
    //      // 400Hz未満なら、低い周波数ほど最大1.5倍までゲインを上げる
    //      // (400 - freq) / 400 で 0.0〜1.0 のブレンド率を作り、0.5fを掛ける
    //      bassBoost += ( (400.0f - freq) / 400.0f ) * 0.1f; 
    //  }*/
      
    //  // 最終的なゲインを適用
    //  finalProfile[h][1] = baseProfile[h][1] * gain;
    //  //finalProfile[h][1] *= bassBoost;
    //}
    
    // 有効なピークの数をカウントするための変数
    int validPeakCount = 0;
    float energy = 0;

    for (int h = 0; h < numPeaks; h++) { 
      // 0.05未満のデータは完全に無視してスキップ
      if(baseProfile[h][1] < ampThreshold) { 
        continue; 
      }
      
      // 有効なデータだけを finalProfile の前から順に詰める
      finalProfile[validPeakCount][0] = baseProfile[h][0]; // 周波数 (Hz) 固定
      finalProfile[validPeakCount][1] = baseProfile[h][1]; // 振幅の初期値
      
      // 除外されなかった「有効なピークだけ」でエネルギーを計算
      energy += baseProfile[h][1] * baseProfile[h][1];
      
      // 次の保存先へインデックスを進める
      validPeakCount++; 
    }
    
    // ベロシティに応じた「目標の全体音量（エネルギーの総量）」の計算
    float targetVolume = basePower * pow(normVel, 0.9f);
    
    // ゲインの算出
    float gain = targetVolume / sqrt(energy + 1e-9f);
    
    // 最終的なゲインを適用（※元の numPeaks ではなく、validPeakCount まで回す！）
    for (int i = 0; i < validPeakCount; i++) {
      // すでに finalProfile に格納されている振幅に対してゲインを掛け算する
      finalProfile[i][1] *= gain; 
    }
    
    return finalProfile;
  }
  
  
  
  
  
  float[][] GetLUT(String type) {
    switch (type) {
      case "Note":
        return note;
      
      case "PianoDecayFF":
        return pianoDecayFFLUT;
      case "PianoWobbleRateFF":
        return pianoWobbleRateFFLUT;
      case "PianoWobbleDepthFF":
        return pianoWobbleDepthFFLUT;
        
      case "PianoDecayMF":
        return pianoDecayMFLUT;
      case "PianoWobbleRateMF":
        return pianoWobbleRateMFLUT;
      case "PianoWobbleDepthMF":
        return pianoWobbleDepthMFLUT;
        
      case "PianoDecayPP":
        return pianoDecayPPLUT;
      case "PianoWobbleRatePP":
        return pianoWobbleRatePPLUT;
      case "PianoWobbleDepthPP":
        return pianoWobbleDepthPPLUT;
        
      case "Kick":
        return kickLUT;
      case "Snare":
        return snareLUT;
      case "HiHat":
        return hihatLUT;
        
      default:
        println("Error: LUT type '" + type + "'が見つかりません!!");
        return null; // 存在しない名前が指定された場合は null を返す
    }
  }
};
\end{lstlisting}

\subsection{InstrumentPlayer.pde}
\begin{lstlisting}[caption={InstrumentPlayer.pde}, label={lst:proc_player}]
class InstrumentPlayer {
  AudioOutput out;

  InstrumentPlayer(AudioOutput out) {
    this.out = out; // Mainから渡された出力を保持する
  }

  void PlayFlute(float midiNote, float duration, float startVel, float endVel, float masterVol) {
    out.playNote(0.0f, duration, 
      new FluteInstrument(midiNote, duration, startVel, endVel, masterVol));
  }
  
  void PlayStrings(float midiNote, float duration, float startVel, float endVel, float masterVol) {
    out.playNote(0.0f, duration, 
      new StringsInstrument(midiNote, duration, startVel, endVel, masterVol));
  }
  
  void PlayPiano(float midiNote, float duration, float velocity, float masterVol) {
    out.playNote(0.0f, duration, 
      new PianoInstrument(midiNote, velocity, masterVol));
  }
   
  // キックを鳴らす関数
  void PlayKick(float masterVol, float velocity) {
    // out.playNote(開始時間, 鳴らす長さ, 音色クラス);
    // ドラムはエンベロープ側で長さを制御するため、第2引数のdurationは適当な値(1.0など)でOKです。
    //out.playNote(0, 0.1, new KickInstrument(masterVol, velocity));
    out.playNote(0, 0.1, new KickInstrument(masterVol, velocity));
  }
  
  // スネアを鳴らす関数
  void PlaySnareOld(float masterVol, float velocity) {
    out.playNote(0, 0.5, new SnareInstrumentOld(masterVol, velocity));
  }
  
  // スネアver2を鳴らす関数
  void PlaySnare(float masterVol, float velocity, float sinAmp) {
    out.playNote(0, 0.3, new SnareInstrument(masterVol, velocity, sinAmp));
  }
  
  // クローズハイハットを鳴らす関数
  void PlayHiHatOld(float masterVol, float velocity) {
    out.playNote(0, 0.2, new HiHatInstrumentOld(masterVol, velocity));
  }
  
  // クローズハイハットver2を鳴らす関数
  void PlayHiHat(float masterVol, float velocity) {
    out.playNote(0, 0.15, new HiHatInstrument(masterVol, velocity));
  }
  
  
  
  
  
  
  
  void PlayHorn(float midiNote, float duration, float startVel, float endVel, float masterVol) {
    out.playNote(0.0f, duration, 
      new HornInstrument(midiNote, duration, startVel, endVel, masterVol));
  }
  void PlayTrumpet(float midiNote, float duration, float startVel, float endVel, float masterVol) {
    out.playNote(0.0f, duration, 
      new TrumpetInstrument(midiNote, duration, startVel, endVel, masterVol));
  }
  
  //void PlayCymbal(float masterVol, float velocity) {
  //  out.playNote(0, 0.2, new CymbalInstrument(masterVol, velocity));
  //}
}
\end{lstlisting}

\subsection{NoteSampler.pde}
\begin{lstlisting}[caption={NoteSampler.pde}, label={lst:proc_sampler}]
// 既存の変数に volatile をつける
volatile int tick;
volatile long interval;
volatile long lastTime = 0;
volatile int[][] returnTick; // {トリガーtick, ジャンプ先tick, ジャンプ先noteIndex}
volatile int returnIndex;
volatile int noteIndex;
volatile boolean isPlaying = false;

// 音楽専用スレッド用の変数を追加
Thread audioThread;




long GetInterval (int playBPM) {
  return 60000000000L / (playBPM * 24);
}



void SoundPlay() {
  if (tick == returnTick[returnIndex][0]) {
    tick = returnTick[returnIndex][1];
    noteIndex = returnTick[returnIndex][2];
    returnIndex ++;
    if (returnIndex == returnTick.length) returnIndex = 0;
  }
  
  if (noteIndex < Note.length) {
    //long a = System.nanoTime();
    float durationTuning = interval / 1000000000.0f;
    while (Note[noteIndex][0] == tick) {
      float duration = float(Note[noteIndex][1]) * durationTuning;
      float pitch = float(Note[noteIndex][2]);
      float startVelocity = float(Note[noteIndex][3]);
      float endVelocity = float(Note[noteIndex][4]);          //startVelocity=127; endVelocity=127;
      float type = (Note[noteIndex].length == 6) ? float(Note[noteIndex][5]) : 0;
      switch (timbre) {
        case "Flute": case "Flute2":
          player.PlayFlute(pitch, duration, startVelocity, endVelocity, 1.0f);
          break;
        case "Strings":
          player.PlayStrings(pitch, duration, startVelocity, endVelocity, 1.0f);
          break;
        case "Piano":
          if      (type == 0) player.PlayPiano(pitch, duration, startVelocity, 1.0f);
          else if (type == 1) player.PlayTrumpet(pitch, duration, startVelocity, endVelocity, 1.0f);
          else if (type == 2) player.PlayHorn(pitch, duration, startVelocity, endVelocity, 1.0f);
          break;
        case "Drum":
          int a = int(pitch) % 12;
          if      (a == 0) player.PlayKick(1.0f, startVelocity);
          //else if (a == 1) player.PlaySnareOld(1.0f, startVelocity);
          //else if (a == 2) player.PlayHiHatOld(1.0f, startVelocity);
          else if (a == 3) player.PlaySnare(1.0f, startVelocity, endVelocity);
          else if (a == 4) player.PlayHiHat(1.0f, startVelocity);
          break;
        case "Horn":
          player.PlayHorn(pitch, duration, startVelocity, endVelocity, 1.0f);
          break;
        case "Trumpet":
          player.PlayTrumpet(pitch, duration, startVelocity, endVelocity, 1.0f);
          break;
        default:
          break;
      }
      noteIndex ++;
      if (noteIndex == Note.length) break;
    }
    //println(System.nanoTime() - a);
  }
}





int[][] Note = {
// ドラム用楽譜（setupでの上書きをコメント文にすれば演奏可能）


  // 【1小節目】
  {0, 96, 72, 96, 96, 0}, // [0] C5

  // 【3小節目】
  {192, 96, 74, 96, 96, 0}, // [1] D5

  // 【5小節目】
  {384, 96, 76, 96, 96, 0}, // [2] E5

  // 【7小節目】
  {576, 96, 77, 96, 96, 0}, // [3] F5

  // 【9小節目】
  {768, 96, 79, 96, 96, 0}, // [4] G5

  // 【11小節目】
  {960, 96, 81, 96, 96, 0}, // [5] A5

  // 【13小節目】
  {1152, 96, 83, 96, 96, 0}, // [6] B5

};





// int[][] Note;
\end{lstlisting}

\subsection{AddEffect.pde}
\begin{lstlisting}[caption={AddEffect.pde}, label={lst:proc_addeffect}]
/*import ddf.minim.AudioEffect;

// 講義の「たたみ込み積分」をそのまま再現した自作エフェクトクラス
class SimpleConvolutionEffect implements AudioEffect {
  float[] impulseResponse; // h(τ) に相当するインパルス応答データ
  float[] historyBuffer;   // 過去の音を記憶しておくバッファ
  int bufferIndex = 0;
  int kernelSize;          // ★過去の音をどれくらい参照するか（サンプル数）

  SimpleConvolutionEffect(int size, float levrl) {
    this.kernelSize = size;
    impulseResponse = new float[kernelSize];
    historyBuffer = new float[kernelSize];
    
    // 簡易的なインパルス応答を自動生成（指数関数的に減衰するノイズ ＝ 擬似的な部屋の響き）
    for (int i = 0; i < kernelSize; i++) {
      float decay = exp(-3.8f * i / kernelSize); // 過去にいくほど音が小さくなる重み
      impulseResponse[i] = random(-1.0f, 1.0f) * decay * levrl; // 響きの強さを調整
    }
  }

  // モノラル信号（現在のシステム）にリアルタイムでたたみ込み演算を行う関数
  void process(float[] signal) {
    for (int i = 0; i < signal.length; i++) {
      float input = signal[i];
      
      // 1. 最新の入力を過去の記憶バッファに格納
      historyBuffer[bufferIndex] = input;
      
      // 2. たたみ込み積分の計算（スライドの数式の再現）
      float output = 0;
      for (int j = 0; j < kernelSize; j++) {
        // 過去へ遡るインデックスを計算
        int idx = (bufferIndex - j + kernelSize) % kernelSize;
        output += historyBuffer[idx] * impulseResponse[j];
      }
      
      // 3. 元のクリーンな音に、生成された反響音を混ぜて出力
      signal[i] = input + output;
      
      // 4. 記憶バッファの書き込み位置を1つ進める
      bufferIndex = (bufferIndex + 1) % kernelSize;
    }
  }

  // ステレオ用（Minimの仕様上必要ですが、中身はモノラル処理を両チャンネルに適用）
  void process(float[] left, float[] right) {
    process(left);
    process(right);
  }
}

*/





/*
// Minimの最新推奨規格「UGen」に従って生まれ変わった反響音クラス
class SimpleConvolutionUGen extends UGen {
  // 前のミキサー（stringsBus）から流れてくる音声信号を受け取るための「ポート（入り口）」
  public UGenInput audio;
  
  float[] impulseResponse; // h(τ) に相当するインパルス応答データ
  float[] historyBuffer;   // 過去の音を記憶しておくバッファ
  int bufferIndex = 0;
  int kernelSize;          // 過去の音をどれくらい参照するか（サンプル数）

  SimpleConvolutionUGen(int size, float level) {
    // 入力ポートを初期化
    audio = new UGenInput(InputType.AUDIO);
    
    this.kernelSize = size;
    impulseResponse = new float[kernelSize];
    historyBuffer = new float[kernelSize];
    
    // 簡易的なインパルス応答を自動生成（指数関数的に減衰する部屋の響き）
    for (int i = 0; i < kernelSize; i++) {
      float decay = exp(-3.8f * i / kernelSize);
      impulseResponse[i] = random(-1.0f, 1.0f) * decay * level;
    }
  }

  // UGenの心臓部：Minimが音声を1サンプル処理するたびに自動で呼び出される関数
  protected void uGenerate(float[] channels) {
    // 1. 直前のミキサー（stringsBus）から流れてきた「現在の入力音」を1サンプル取得
    float input = audio.getLastValue();
    
    // 2. 最新の入力を過去の記憶バッファに格納
    historyBuffer[bufferIndex] = input;
    
    // 3. たたみ込み積分の計算（講義スライドの数式の完全再現）
    float output = 0;
    for (int j = 0; j < kernelSize; j++) {
      int idx = (bufferIndex - j + kernelSize) % kernelSize;
      output += historyBuffer[idx] * impulseResponse[j];
    }
    
    // 4. 記憶バッファの書き込み位置を1つ進める
    bufferIndex = (bufferIndex + 1) % kernelSize;
    
    // 5. 元の音に反響音を足し合わせる
    float finalOutput = input + output;
    
    // 6. 出力先へ音を流し込む（ループにすることでモノラル・ステレオどちらにも自動対応！）
    for (int i = 0; i < channels.length; i++) {
      channels[i] = finalOutput;
    }
  }
}*/






// Minimの最新推奨規格「UGen」に従って生まれ変わった反響音クラス（完全ステレオ対応版）
class AddReverb extends UGen {
  public UGenInput audio;
  
  float[] impulseResponse;
  
  // ★改造ポイント1：記憶バッファを左右の耳で完全に独立させる
  float[] historyBufferL;
  float[] historyBufferR;
  
  // UGenの uGenerate は1サンプルフレーム（左右同時）ごとに呼ばれるため、時間は共通（1つ）でOK
  int bufferIndex = 0;
  int kernelSize;

  AddReverb(int size, float level) {
    audio = new UGenInput(InputType.AUDIO);
    this.kernelSize = size;
    
    impulseResponse = new float[kernelSize];
    historyBufferL = new float[kernelSize];
    historyBufferR = new float[kernelSize];
    
    // インパルス応答の生成
    for (int i = 0; i < kernelSize; i++) {
      float decay = exp(-3.8f * i / kernelSize);
      impulseResponse[i] = random(-1.0f, 1.0f) * decay * level;
    }
  }

  // UGenの心臓部
  protected void uGenerate(float[] channels) {
    // ★改造ポイント2：入力音を「配列」として取得し、左右の音を別々に取り出す
    float[] inputs = audio.getLastValues();
    float inL = inputs[0]; 
    // もし入力がモノラル（1ch）だった場合は、右耳用にも左と同じ音を入れるセーフティガード
    float inR = (inputs.length > 1) ? inputs[1] : inL;
    
    // 最新の入力を左右それぞれの記憶バッファに格納
    historyBufferL[bufferIndex] = inL;
    historyBufferR[bufferIndex] = inR;
    
    // ★改造ポイント3：たたみ込み積分の計算を、左右独立して行う
    float outL = 0;
    float outR = 0;
    for (int j = 0; j < kernelSize; j++) {
      int idx = (bufferIndex - j + kernelSize) % kernelSize;
      outL += historyBufferL[idx] * impulseResponse[j];
      outR += historyBufferR[idx] * impulseResponse[j];
    }
    
    // 記憶バッファの書き込み位置を1つ進める（左右共通の時間が進む）
    bufferIndex = (bufferIndex + 1) % kernelSize;
    
    // ★改造ポイント4：出力先（channels）がモノラルかステレオかで安全に音を振り分ける
    if (channels.length == 1) {
      // 出力先がモノラルの場合
      channels[0] = inL + outL;
    } else if (channels.length >= 2) {
      // 出力先がステレオの場合
      channels[0] = inL + outL; // 左チャンネルの出力
      channels[1] = inR + outR; // 右チャンネルの出力
    }
  }
}
\end{lstlisting}

\subsection{Flute.pde}
\begin{lstlisting}[caption={Flute.pde}, label={lst:proc_flute}]
// ==========================================
// トリル専用 LFOウェーブテーブル（台形波）の生成と保持
// ==========================================
Wavetable trillWave = CreateTrapezoidWave();

// 台形波（角の丸いパルス波）を生成して返す関数
Wavetable CreateTrapezoidWave() {
  int waveSize = 1024; // 波形の解像度
  float[] waveData = new float[waveSize];
  
  // 傾斜（スロープ）の幅を設定。
  // 全体の何%を「移動時間」に使うか。例えば 0.15 なら 15% がスロープ、85%が平らな部分になる。
  float slopeRatio = 0.5f;
  // この数値を 0.05（より矩形波に近い、素早い指の動き）や 0.3（もったりとしたゆっくりな指の動き）に変更することで
  // トリルの「ニュアンス（奏者の癖）」までプログラミングできるようになる。
  
  for (int i = 0; i < waveSize; i++) {
    float phase = (float)i / waveSize; // 1周期を -1.0 〜 +1.0 の間で描く
    
    if (phase < 0.5f) { // 前半（高い音に向かう部分）
      float localPhase = phase / 0.5f; // 0.0 〜 1.0 に正規化
      if (localPhase < slopeRatio) {
        // 立ち上がりのスロープ（-1.0 から +1.0 へ滑らかに移動）
        // smoothstep関数(3x^2 - 2x^3)を使って、S字カーブを描いてより滑らかにする
        float t = localPhase / slopeRatio;
        waveData[i] = -1.0f + 2.0f * (3*t*t - 2*t*t*t); // smoothstep
      } else {
        // 平らな部分（高い音を維持）
        waveData[i] = 1.0f;
      }
    
    } else {
      // 後半（低い音に戻る部分）
      float localPhase = (phase - 0.5f) / 0.5f; // 0.0 〜 1.0 に正規化
      if (localPhase < slopeRatio) {
        // 立ち下がりのスロープ（+1.0 から -1.0 へ滑らかに移動）
        float t = localPhase / slopeRatio;
        waveData[i] = 1.0f - 2.0f * (3*t*t - 2*t*t*t); // smoothstep
      } else {
        // 平らな部分（低い音を維持）
        waveData[i] = -1.0f;
      }
    }
  }
  
  return new Wavetable(waveData);
}












class FluteInstrument implements Instrument {
  byte numHarmonics/* = 15*/;
  
  ADSR masterEnv;
  ADSR[] harmonicEnvs/* = new ADSR[numHarmonics]*/;
  
  Line[] wobbleLines;    // 6/12【追加】倍音ウナリ（AM）用のフェードライン
  float[] wobbleFactors; // 6/12【追加】倍音ごとの固有ウナリ比率の保持用
  
  Line[] ampLines/* = new Line[numHarmonics]*/;
  float[] startAmps/* = new float[numHarmonics]*/;
  float[] endAmps/* = new float[numHarmonics]*/;
  
  ADSR breathEnv;
  
  
  // ==========================================
  // 【新規追加】 持続する空気の渦と管の共鳴（サブハーモニクス層）
  // ==========================================
  ADSR sustainEnv;
  
  // 【追加】持続ノイズの音量変化（クレッシェンド/デクレッシェンド）用
  Line sustainAmpLine;
  Line sustainWobbleLine; // 6/12【追加】持続ノイズのウナリ用のフェードライン
  float startSustainVol;
  float endSustainVol;
  
  
  
  FluteInstrument(float midiNote, float duration, float startVel, float endVel, float masterVol) {
    Summer mixer = new Summer();
    
    //float masterAttackTime = 0.2f - constrain(startVel, 0, 127) / 127.0f * 0.15f;
    float masterAttackTime = 0.18f;
    masterAttackTime = min(masterAttackTime, duration - 0.1f);
    if (masterAttackTime < 0.035f) masterAttackTime = 0.035f;
    
    // 【修正】masterEnv の一番最後の引数（リリース時間）を 0.3f から 0.4f に延ばす。
    // 内部の音が 0.3秒 かけて完全に消え去るのを待ってから、安全にアンパッチさせるための余裕（バッファ）です。
    // 【修正】masterEnv のリリース時間を、希望のフェードアウト時間（例：0.35秒）に設定
    masterEnv = new ADSR(1.0f, masterAttackTime, 0.1f, 0.9f, 0.2f);
    
    mixer.patch(masterEnv);
    
    // ==========================================
    // 【新規】トリル判定と実音の計算
    // ==========================================
    boolean isTrill = midiNote >= 128;
    float actualMidiNote = isTrill ? midiNote - 128 : midiNote;
    
    // 実音の周波数と、トリル先（全音上＝+2）の周波数
    //float f0 = Frequency.ofMidiNote(actualMidiNote).asHz();
    //float f0_trill = Frequency.ofMidiNote(actualMidiNote + 1).asHz(); // 半音トリルにしたい場合は +1 にする

    // LUTの倍音プロファイル
    // 10倍音、最低音59、LUT配列、フルートのパワー(1.4f)
    //float[][] startProfile = utils.GetDynamicProfile( actualMidiNote, startVel, "Flute", 59, 1.3f);
    //float[][] endProfile   = utils.GetDynamicProfile( actualMidiNote, endVel, "Flute", 59, 1.3f);
    
    //float volFactor = masterVol * 0.11f; 
    
    
    //numHarmonics = (byte)startProfile.length;
    // 【重要：NPE解決策】実際のデータ数に合わせて、動的に配列のメモリを確保する！
    //harmonicEnvs = new ADSR[numHarmonics];
    //ampLines     = new Line[numHarmonics];
    //startAmps    = new float[numHarmonics];
    //endAmps      = new float[numHarmonics];
    
    
    // 変更6/12
    // 1. 倍音レシピを安全に取り出すための基準ベロシティ（0除算・ロード不全の防止）
    float profileVel = (startVel > 0) ? startVel : endVel;
    if (profileVel <= 0) profileVel = 64.0f; // どちらも0の場合のセーフティ
    
    // 2. 開始・終了のゲインファクターをフルートの特性（0.7乗カーブ）に基づいて独立計算
    float startGainFactor = pow(startVel / profileVel, 0.9f);
    float endGainFactor   = pow(endVel / profileVel, 0.9f);
    
    // 安全にロードされた基準プロファイルを取得
    float[][] startProfile = utils.GetDynamicProfile(actualMidiNote, profileVel, "Flute", 59, 1.3f);
    
    float volFactor = masterVol * 0.022f;
    
    numHarmonics = (byte)startProfile.length;
    
    harmonicEnvs  = new ADSR[numHarmonics];
    ampLines      = new Line[numHarmonics];
    wobbleLines   = new Line[numHarmonics];    // 初期化を追加
    wobbleFactors = new float[numHarmonics];   // 初期化を追加
    startAmps     = new float[numHarmonics];
    endAmps       = new float[numHarmonics];
    
    
    
    
    //float fundamentalFreq = 440.0 * pow(2.0, (actualMidiNote - 69) / 12.0);
    byte baseWaveNum = 0;
    for (byte i = 1; i < numHarmonics; i++) {
      if (startProfile[i-1][0] < startProfile[i][0]) baseWaveNum = i;
    }
    
    
    // ==========================================
    // 【新規】倍音ごとの自然な振幅揺らぎ（変動率ベース）
    // ==========================================
    // Pythonで解析した変動率(%)を参考に、基音の振幅に対して何倍の揺れ(±)を許容するかを設定。
    // 実測データの変動率(%)から数式 [ Variation / 500 ] によって導出された、
    // 倍音ごとの自然な振幅揺らぎ（AM変調）の深度設定。
    float[] wobbleDepths = {0.1250f, 0.1616f, 0.2078f, 0.2340f, 0.2114f, 0.2742f, 0.3716f, 0.2946f, 0.7420f, 0.8796f, 
                            0.7850f, 0.5070f, 0.6242f, 0.7654f, 0.6702f, 0.3898f, 0.8850f, 0.5942f, 0.3152f, 0.5030f};
    
    // ==========================================
    // 【新規】この音符全体の「呼吸のスピード」を決定
    // forループの外で1つだけ定義し、全倍音で共有することで相殺を防ぐ
    // ==========================================
    float globalBreathSpeed = random(2.5f, 3.5f); // 生楽器の自然なビブラート周期
    
    
    


    // 1. 加算合成パート（7つのサイン波）
    for (int i = 0; i < numHarmonics; i++) {
      /*
      int h = i + 1;
      // あなたがPythonで導き出した厳密なインハーモニシティ公式を適用
      float B = 0.0005f; // インハーモニシティ係数
      float inharmonicity = sqrt(1.0f + B * h * h);
      float targetFreq = f0 * h * inharmonicity;
      */
      float targetFreq = startProfile[i][0];
      
      Oscil osc = new Oscil(targetFreq, 0.0f, Waves.SINE);
      
      // 【追加】初期位相を 0.0 〜 1.0 (0度〜360度) の間で完全にランダムに散らす
      // これにより「カチッ」とした機械的なアタック音が消え、音割れ（ピーク）も防げる
      osc.setPhase(random(1.0f));
      
      
      //startAmps[i] = startProfile[i][1] * volFactor;
      //endAmps[i] = endProfile[i] * volFactor;
        // 聴覚特性（0.7乗）に完全同期させた、開始時から終了時への音量変化比率を算出
        //float velRatio = pow(endVel / (startVel + 0.0001f), 0.7f);
      // 【確定修正】すでにvolFactorが掛けられたstartAmpsに対して、正しいべき乗比率を乗算する
      //endAmps[i]   = startAmps[i] * pow(endVel / (startVel + 0.0001f), 0.7f);
      
      // 6/12追加
      // 独立した開始・終了音量を厳密に算出
      float baseAmp = startProfile[i][1] * volFactor;
      startAmps[i] = baseAmp * startGainFactor;
      endAmps[i]   = baseAmp * endGainFactor;
      
      
      // ==========================================
      // 【修正】LFOをトリルとビブラートで分岐
      // ==========================================
      //if (isTrill) {
      //  // トリル時の上限周波数
      //  //float trillFreq = f0_trill * (i+1) * inharmonicity;
      //  float trillFreq = targetFreq * pow(2.0f, 1.0f/12.0f);
        
      //  // 矩形波LFOの中心を2つの音の真ん中に置き、振幅を「差の半分」にする
      //  float centerFreq = (targetFreq + trillFreq) / 2.0f;
      //  float ampFreq = abs(trillFreq - targetFreq) / 2.0f;
        
      //  // 7.5Hzのスピードで生成したグローバル変数 trillWave を使って瞬間的に音を切り替える
      //  Oscil freqController = new Oscil(6.5f, ampFreq, trillWave);
      //  freqController.offset.setLastValue(centerFreq);
        
      //  // トリル時のみ、オシレーターの周波数にLFOを接続する
      //  freqController.patch(osc.frequency);
        
      //  wobbleLines[i] = null;   // トリル時はAM不要のため安全ガード
      //  wobbleFactors[i] = 0.0f;
        
      //} else {
        // 通常のビブラート（サイン波）
        
        // ==========================================
        // ① 極小のピッチ揺れ (FM変調) の復活
        // 以前の0.0132fを「0.003f」に激減させ、音痴にならない程度の「空気の揺らぎ」を作る
        // ==========================================
        float vibDepthHz = targetFreq * 0.0015f; 
        Oscil freqController = new Oscil(globalBreathSpeed, vibDepthHz/10, Waves.SINE);
        freqController.offset.setLastValue(targetFreq); 
        freqController.patch(osc.frequency);
        
        // ==========================================
        // ② 呼吸に同期した音量揺らぎ (AM変調)
        // ==========================================
        // この倍音の基準音量に対して、配列で設定した割合だけ揺らす
        //float wobbleAmp = startAmps[i] * wobbleDepths[i]; 
        //float wobbleAmp = startAmps[i] * random(0.1f, 0.3f);
        
        // 1.5Hz 〜 3.5Hz の間で、倍音ごとにランダムなスピードで揺らす
        // 【修正】ランダムではなく、統一された globalBreathSpeed を使う！
        //Oscil ampLFO = new Oscil(globalBreathSpeed, wobbleAmp, Waves.SINE);
        
        // 位相（スタート位置）は少しだけバラけさせると、音が立体的になります
        //ampLFO.setPhase(random(1.0f)); 
        
        // ampLFOを osc.amplitude にパッチ(追加)する。
        // Minimでは複数のUGenを同じ入力にパッチすると「加算(Sum)」されるため、
        // ampLines の基準値に対して、ampLFO が ±wobbleAmp の揺れを足し引きしてくれます。
        //ampLFO.patch(osc.amplitude);
        
        // --- 音量揺らぎ（AM変調）のフェード構造化 ---
        wobbleFactors[i] = random(0.1f, 0.2f); 
        wobbleFactors[i] = wobbleDepths[i]; 
        wobbleLines[i] = new Line();
        Oscil ampLFO = new Oscil(globalBreathSpeed, 0.0f, Waves.SINE); // 初期振幅は0
        ampLFO.setPhase(random(1.0f)); 
        
        wobbleLines[i].patch(ampLFO.amplitude); // 新設LineをLFOの振幅へ接続
        ampLFO.patch(osc.amplitude);
      //}
      
      ampLines[i] = new Line();
      ampLines[i].patch(osc.amplitude); 
      
      float attackTime = (i == baseWaveNum) ? 0.058f : 0.395f; 
      
      // 【修正】内部のサイン波のリリースを 0.3f から 10.0f（10秒）にして、急激な角を作らせない
      harmonicEnvs[i] = new ADSR(1.0f, attackTime/3, 1.0f, 0.7f, 10.0f);
      osc.patch(harmonicEnvs[i]);
      harmonicEnvs[i].patch(mixer);
    }
    
    
    
    // ==========================================
    // 2. 息のノイズ（Chiff）パートの実測値アップデート
    // =========================================
    // 0.0〜1.0 のベロシティ割合を計算（ノイズの音量調整用）
    float startNorm = constrain(startVel, 0, 127) / 127.0f;
    
    // 中低音の破裂音（ドスッという音）になるため、少しだけ音量を上げます（0.01f -> 0.02f）
    float attackNoiseVol = masterVol * 0.008f * pow(startNorm, 0.9f); 
    Noise breathNoise = new Noise(attackNoiseVol, Noise.Tint.WHITE);
    
    // 【修正】ハイパス(HP)からバンドパス(BP)に変更。
    // 3500Hz付近の「シュッ」という音だけを残し、耳障りな高音を削る。
    // 【修正】実測データに基づき、3500Hzから 581.4Hz に変更！
    // 473Hzの山も一緒に拾うため、レゾナンスを 0.3 から 0.15 に下げて帯域を広げる
    MoogFilter breathFilter = new MoogFilter(581.4f, 0.15f);
    breathFilter.type = MoogFilter.Type.BP;
    
    breathEnv = new ADSR(1.0f, 0.001f, 0.2f, 0.0f, 0.0f);
    
    breathNoise.patch(breathFilter);
    breathFilter.patch(breathEnv);
    breathEnv.patch(mixer);
    
    
    /*
    // ==========================================
    // 3. 【新規追加】持続する空気の渦と管の共鳴（サブハーモニクス層）
    // ==========================================
    float endNorm = constrain(endVel, 0, 127) / 127.0f; // 【追加】終了時のベロシティ割合
    
    // 【修正】開始時と終了時のノイズ音量を計算して保持する
    startSustainVol = masterVol * 0.005f * pow(startNorm, 0.9f);
    endSustainVol   = masterVol * 0.005f * pow(endNorm, 0.9f);
    
    // ピンクノイズで低音の「フオォォ」という空気感を出す
    // 【修正】Noiseの初期音量は0.0fにし、音量制御はLineに任せる
    Noise sustainNoise = new Noise(0.0f, Noise.Tint.PINK);
    
    // 【追加】持続ノイズのボリュームつまみにLineを接続
    sustainAmpLine = new Line();
    sustainAmpLine.patch(sustainNoise.amplitude);
    
    // ==========================================
    // 【新規】持続ノイズ自体にも「息の乱れ」を適用する
    // サイン波の第1倍音と同等のスピードと深さでノイズを揺らす
    // ==========================================
    //float noiseWobbleAmp = startSustainVol * 0.20f; // 20%揺らす
    //Oscil sustainWobble = new Oscil(random(1.5f, 3.0f), noiseWobbleAmp, Waves.SINE);
    //sustainWobble.setPhase(random(1.0f));
    //sustainWobble.patch(sustainNoise.amplitude); // Lineの音量にLFOの揺れを加算
    // ==========================================
    
    // 6/12変更
    // --- 持続ノイズの息の乱れも完全にフェードさせる ---
    sustainWobbleLine = new Line();
    Oscil sustainWobble = new Oscil(random(1.5f, 3.0f), 0.0f, Waves.SINE); // 初期振幅は0
    sustainWobble.setPhase(random(1.0f));
    
    sustainWobbleLine.patch(sustainWobble.amplitude); // 新設LineをノイズLFOの振幅へ接続
    sustainWobble.patch(sustainNoise.amplitude);      // LFO出力をノイズ振幅へ（加算）
    
    // ローパスフィルターで高音を削りつつ、2500Hz付近に共鳴(レゾナンス0.5)の山を作る
    MoogFilter sustainFilter = new MoogFilter(2500f, 0.5f);
    sustainFilter.type = MoogFilter.Type.LP;
    
    // 音が鳴っている間ずっと持続するエンベロープ（サイン波のアタックに合わせる）
    // 【修正】持続ノイズのリリースも 0.3f から 10.0f（10秒）にする
    sustainEnv = new ADSR(1.0f, 0.058f, 0.0f, 1.0f, 10.0f);
    
    sustainNoise.patch(sustainFilter);
    sustainFilter.patch(sustainEnv);
    sustainEnv.patch(mixer);
    */
  }
  
  
  /*
  void noteOn(float dur) {
    masterEnv.noteOn();
    
    // ====================================================
    // 【修正】Lineの寿命に、リリース時間（0.4秒）分の猶予を足す！
    // これにより、減衰中もLineが生き残り、オシレーターの音量が突然0.0fにリセットされるのを防ぐ
    // ====================================================
    // 【微調整】masterEnvのリリース時間(0.35f)より少し長ければOK
    float safeDur = dur + 0.5f;
    
    for(int i=0; i<numHarmonics; i++) {
      harmonicEnvs[i].noteOn();
      ampLines[i].activate(safeDur, startAmps[i], endAmps[i]);
    }
    
    breathEnv.noteOn();
    sustainEnv.noteOn(); // 【追加】持続ノイズをオン
    sustainAmpLine.activate(safeDur, startSustainVol, endSustainVol); // 【追加】発音と同時に、持続ノイズの音量変化をスタートさせる
    
    masterEnv.patch(out); 
  }*/
  // 6/12変更
  void noteOn(float dur) {
    masterEnv.noteOn();
    float safeDur = dur + 0.5f;
    
    for(int i = 0; i < numHarmonics; i++) {
      harmonicEnvs[i].noteOn();
      
      // 1. 各倍音のベース音量をフェード
      ampLines[i].activate(safeDur, startAmps[i], endAmps[i]);
      
      // 2. 通常ビブラート時のみ、ウナリ（AM）の深さも比例フェードさせる
      if (wobbleLines[i] != null) {
        float factor = wobbleFactors[i];
        wobbleLines[i].activate(safeDur, startAmps[i] * factor, endAmps[i] * factor);
      }
    }
    
    breathEnv.noteOn();
    //sustainEnv.noteOn(); 
    
    // 3. 持続ノイズの全体音量と、そのウナリの深さを同時に追従フェード
    //sustainAmpLine.activate(safeDur, startSustainVol, endSustainVol); 
    //sustainWobbleLine.activate(safeDur, startSustainVol * 0.20f, endSustainVol * 0.20f); 
    
    masterEnv.patch(longOut); 
  }
  
  
  
  void noteOff() {
    masterEnv.noteOff();
    for(int i=0; i<numHarmonics; i++) harmonicEnvs[i].noteOff();
    
    // 【削除】breathEnv は既に無音なので noteOff を呼ばない（空撃ちグリッチ防止）
    // breathEnv.noteOff();
    
    //sustainEnv.noteOff(); // 【追加】持続ノイズをオフ
    masterEnv.unpatchAfterRelease(longOut); 
  }
}
\end{lstlisting}

\subsection{Strings.pde}
\begin{lstlisting}[caption={Strings.pde}, label={lst:proc_strings}]
import java.util.Arrays;
import java.util.Comparator;

// ==========================================
// 究極のストリングス・インストゥルメント (ビブラート数式モデル実装版)
// ==========================================
class StringsInstrument implements Instrument {
  //byte numHarmonics = 30;
  //Summer mixer;
  ADSR masterEnv; 
  //ADSR[] harmonicEnvs = new ADSR[numHarmonics];
  
  /*
  Line[] ampLines = new Line[numHarmonics];
  float[] startAmps = new float[numHarmonics];
  float[] endAmps = new float[numHarmonics];
  
  // ビブラートの深さをコントロールするためのLine
  Line[] vibDepthLines = new Line[numHarmonics];
  float[] targetVibDepths = new float[numHarmonics]; // 各倍音の最終的な揺れ幅(Hz)を保存
  */
  
  // クラス上部のメンバー変数宣言部分の修正
  byte maxTotalOscils = 40; // クロスフェード時に2つの楽器が混ざるため、多めに確保
  Oscil[] oscils = new Oscil[maxTotalOscils];
  Line[] ampLines = new Line[maxTotalOscils];
  Line[] vibDepthLines = new Line[maxTotalOscils];
  
  float[] startAmps = new float[maxTotalOscils];
  float[] endAmps = new float[maxTotalOscils];
  float[] targetVibDepths = new float[maxTotalOscils];
  
  byte totalActiveOscils = 0; // 今回の音で実際に生成されたオシレーターの総数
  
  Line[] wobbleLines = new Line[maxTotalOscils]; // 【新規追加】ウナリのフェード用ライン 6/12
  
  //Noise biteNoise;
  //MoogFilter biteFilter;
  //ADSR biteEnv;
  //ADSR hissEnv;
  
  //Noise hissNoise;
  //MoogFilter hissFilter;
  //Line hissAmpLine;
  //Oscil hissWobble;
  //float startHissVol, endHissVol;
  
  StringsInstrument(float midiNote, float duration, float startVel, float endVel, float masterVol) {
    Summer mixer = new Summer();
    
    // =========================================================
    // 6/13【新設計】音符の長さに応じた動的アタックエンベロープ
    // =========================================================
    float attackTime = 0.3f - constrain(startVel, 0, 127) / 127.0f * 0.25f;
    attackTime = min(attackTime, duration - 0.1f);
    if (attackTime < 0.035f) attackTime = 0.035f;
    
    
    masterEnv = new ADSR(1.0f, attackTime, 0.5f, 0.9f, 0.12f);
    mixer.patch(masterEnv); 
    
    //float f0 = Frequency.ofMidiNote(midiNote).asHz();
    float volFactor = masterVol * 0.020f; 
    
    //float startNorm = constrain(startVel, 0, 127) / 127.0f;
    //float endNorm   = constrain(endVel, 0, 127) / 127.0f;
    
    
    // =========================================================
    // 【データ分析から導いた数式①】ビブラートスピード
    // 低音(G3=55)はゆっくり、高音(E7=100)は速く。さらに人間らしいランダムな揺らぎを付加。
    // =========================================================
    float baseVibSpeed = map(midiNote, 48, 100, 5.2f, 6.5f);
    float globalVibSpeed = baseVibSpeed + random(-0.2f, 0.2f);
    
    // =========================================================
    // 【データ分析から導いた数式②】ビブラートの深さ（Cents）
    // 低音は浅く(±10cents)、高音は深く(±25cents)揺れる。異常な暴走はここでカットされる。
    // =========================================================
    float fmDepthCents = map(midiNote, 48, 100, 10.0f, 20.0f) * random(0.9f, 1.0f) * 0.2f;
    
    
    /*
    // =========================================================
    // 【クロスフェード処理】オーケストラ・ストリングスの4Wayブレンド
    // =========================================================
    
    // コントラバス：最低音 C1 (MIDI 24)
    float[][] bassStart   = utils.GetDynamicProfile(midiNote, startVel, "Bass", 24, 0.9f);
    float[][] bassEnd     = utils.GetDynamicProfile(midiNote, endVel, "Bass", 24, 0.9f);

    // チェロ：最低音 C2 (MIDI 36)
    float[][] celloStart  = utils.GetDynamicProfile(midiNote, startVel, "Cello", 36, 0.9f);
    float[][] celloEnd    = utils.GetDynamicProfile(midiNote, endVel, "Cello", 36, 0.9f);
    
    // ビオラ：最低音 C3 (MIDI 48)
    float[][] violaStart  = utils.GetDynamicProfile(midiNote, startVel, "Viola", 48, 0.9f);
    float[][] violaEnd    = utils.GetDynamicProfile(midiNote, endVel, "Viola", 48, 0.9f);
    
    // バイオリン：最低音 G3 (MIDI 55)
    float[][] violinStart = utils.GetDynamicProfile(midiNote, startVel, "Violin", 55, 0.9f);
    float[][] violinEnd   = utils.GetDynamicProfile(midiNote, endVel, "Violin", 55, 0.9f);
    
    float[] startProfile = new float[numHarmonics];
    float[] endProfile = new float[numHarmonics];
    
    // ---------------------------------------------------------
    // ベースとなるブレンド割合の計算 (合計が必ず1.0になるように分配)
    // ---------------------------------------------------------
    float baseBassRatio = 0.0f;
    float baseCelloRatio = 0.0f;
    float baseViolaRatio = 0.0f;
    float baseViolinRatio = 0.0f;
    
    if (midiNote <= 36) {
      // C2(36) 以下はコントラバスの独壇場
      baseBassRatio = 1.0f;
    } else if (midiNote <= 48) {
      // C2(36) 〜 C3(48) : コントラバスからチェロへ滑らかにクロスフェード
      baseBassRatio = map(midiNote, 36, 48, 1.0f, 0.0f);
      baseCelloRatio = 1.0f - baseBassRatio;
    } else if (midiNote <= 60) {
      // C3(48) 〜 C4(60) : チェロからビオラへ滑らかにクロスフェード
      baseCelloRatio = map(midiNote, 48, 60, 1.0f, 0.0f);
      baseViolaRatio = 1.0f - baseCelloRatio;
    } else if (midiNote <= 72) {
      // C4(60) 〜 C5(d) : ビオラからバイオリンへ滑らかにクロスフェード
      baseViolaRatio = map(midiNote, 60, 72, 1.0f, 0.0f);
      baseViolinRatio = 1.0f - baseViolaRatio;
    } else {
      // C5(72) 以上はバイオリンの独壇場
      baseViolinRatio = 1.0f;
    }
    
    float bassBlendStart = baseBassRatio, bassBlendEnd = baseBassRatio;
    float celloBlendStart = baseCelloRatio, celloBlendEnd = baseCelloRatio;
    float violaBlendStart = baseViolaRatio, violaBlendEnd = baseViolaRatio;
    float violinBlendStart = baseViolinRatio, violinBlendEnd = baseViolinRatio;
    
    // ---------------------------------------------------------
    // 【特例パッチ】 D6(86) と D#6/Eb6(87) の波形破綻への対応
    // mf(80)まではバイオリン本来の音色。そこからff(112)に向かってビオラを混ぜる
    // ---------------------------------------------------------
    if (midiNote == 86 || midiNote == 87) {
      float threshMF = 80.0f / 127.0f;
      float threshFF = 112.0f / 127.0f;
      
      if (startNorm > threshMF) {
        float clampedStart = constrain(startNorm, threshMF, threshFF);
        violinBlendStart = map(clampedStart, threshMF, threshFF, 1.0f, 0.3f);
        violaBlendStart = 1.0f - violinBlendStart;
      }
      if (endNorm > threshMF) {
        float clampedEnd = constrain(endNorm, threshMF, threshFF);
        violinBlendEnd = map(clampedEnd, threshMF, threshFF, 1.0f, 0.3f);
        violaBlendEnd = 1.0f - violinBlendEnd;
      }
      
      bassBlendStart = 0.0f; bassBlendEnd = 0.0f;
      celloBlendStart = 0.0f; celloBlendEnd = 0.0f;
    }
    // ---------------------------------------------------------
    
    // 4つの楽器のプロファイルを、計算した割合で混ぜ合わせる
    for (int i = 0; i < numHarmonics; i++) {
      startProfile[i] = (bassStart[i] * bassBlendStart) + (celloStart[i] * celloBlendStart) + (violaStart[i] * violaBlendStart) + (violinStart[i] * violinBlendStart);
      endProfile[i]   = (bassEnd[i]   * bassBlendEnd)   + (celloEnd[i]   * celloBlendEnd)   + (violaEnd[i]   * violaBlendEnd)   + (violinEnd[i]   * violinBlendEnd);
    }
    // =========================================================
    
    
    
    // ==================================================
    // 【インハーモニシティ係数(B)の動的計算セクション】
    // ==================================================
    
    // 1. 各楽器の基準B係数（解析結果の中央値から設定）
    float bBassBase   = 0.000070f;
    float bCelloBase  = 0.000065f;
    float bViolaBase  = 0.000055f;
    float bViolinBase = 0.000052f;
    
    // 2. 各楽器の基準となる開放弦周波数（正規化用）
    float fRefBass   = 32.7f;  // C1相当
    float fRefCello  = 65.4f;  // C2相当
    float fRefViola  = 130.8f; // C3相当
    float fRefViolin = 196.0f; // G3相当
    
    // 3. 現在の楽器ブレンド状態に合わせた「基準B」と「基準周波数」を算出
    float blendedBBase = (bBassBase * baseBassRatio) + (bCelloBase * baseCelloRatio) 
                       + (bViolaBase * baseViolaRatio) + (bViolinBase * baseViolinRatio);
                       
    float blendedFRef  = (fRefBass * baseBassRatio) + (fRefCello * baseCelloRatio) 
                       + (fRefViola * baseViolaRatio) + (fRefViolin * baseViolinRatio);
    
    // 4. 【核心】ピッチの2乗に比例してBを増幅させる（ハイポジション効果の再現）
    // あなたの観察通り、C5(523Hz)において約0.00006になるようスケーリングされます
    float B = blendedBBase * pow(f0 / blendedFRef, 2.0f);
    
    // 異常値ガード（物理的に破綻しない範囲に制限）
    B = constrain(B, 0.00001f, 0.0005f);
    
    // ==================================================
    
    
    
    for (int i = 0; i < numHarmonics; i++) {
      int h = i + 1;
      // 動的に計算された B を適用
      float inharmonicity = sqrt(1.0f + B * h * h);
      float targetFreq = f0 * h * inharmonicity; 
      
      // 【改善】アンチエイリアス (20000Hz超えの倍音は、生成処理自体をストップしてCPUを節約)
      if (targetFreq > 22050.0f) {
        continue; // ここでループを抜け、これ以上高い倍音は作らない
      }
      
      float finalStartAmp = startProfile[i] * volFactor;
      float finalEndAmp = endProfile[i] * volFactor;
      
      // アンチエイリアス (20000Hz超えの強制ミュート)
      //if (targetFreq > 20000.0f) {
      //  finalStartAmp = 0.0f;
      //  finalEndAmp = 0.0f;
      //}
      
      Oscil osc = new Oscil(targetFreq, 0.0f, Waves.SINE);
      osc.setPhase(0.0f); // ストリングスの核（エッジ）
      
      // =========================================================
      // ピッチ揺れ (Delayed FM変調) の準備
      // セント(Cents)を周波数(Hz)の揺れ幅に変換して配列に保存しておく
      // =========================================================
      float vibDepthRatio = pow(2.0f, fmDepthCents / 1200.0f) - 1.0f;
      targetVibDepths[i] = targetFreq * vibDepthRatio; 
      
      // 揺れの深さをコントロールするLFO（最初は振幅0でスタート）
      Oscil freqController = new Oscil(globalVibSpeed, 0.0f, Waves.SINE);
      freqController.offset.setLastValue(targetFreq); 
      freqController.patch(osc.frequency);
      
      // ビブラートの深さを0から目標値へ変化させるLine
      vibDepthLines[i] = new Line();
      vibDepthLines[i].patch(freqController.amplitude);
      
      // =========================================================
      // 音量揺らぎ (AM変調) の適用
      // 高次倍音ほど激しく揺れる（15%〜60%）
      // =========================================================
      startAmps[i] = finalStartAmp;
      endAmps[i] = finalEndAmp;
      
      float wobbleDepth = map(i, 0, numHarmonics - 1, 0.15f, 0.60f);
      float wobbleAmp = startAmps[i] * wobbleDepth; 
      
      Oscil ampLFO = new Oscil(globalVibSpeed, wobbleAmp, Waves.SINE);
      ampLFO.setPhase(random(1.0f)); 
      ampLFO.patch(osc.amplitude);
      
      ampLines[i] = new Line();
      ampLines[i].patch(osc.amplitude); 
      
      //harmonicEnvs[i] = new ADSR(1.0f, 0.2f, 0.0f, 1.0f, 0.3f);
      //osc.patch(harmonicEnvs[i]);
      //harmonicEnvs[i].patch(mixer);
      
      // 【修正1】harmonicEnvsを削除し、純粋にミキサーへ送る
      osc.patch(mixer);
    }
    */
    
    // =========================================================
    // 【クロスフェード処理】オーケストラ・ストリングスの4Wayブレンド
    // =========================================================
    // 1. 各楽器のブレンド割合の計算 (合計が必ず1.0になるように分配)
    float baseBassRatio = 0.0f;
    float baseCelloRatio = 0.0f;
    float baseViolaRatio = 0.0f;
    float baseViolinRatio = 0.0f;
    
    if (midiNote <= 36) {
      baseBassRatio = 1.0f;
    } else if (midiNote <= 48) {
      baseBassRatio = map(midiNote, 36, 48, 1.0f, 0.0f);
      baseCelloRatio = 1.0f - baseBassRatio;
    } else if (midiNote <= 60) {
      baseCelloRatio = map(midiNote, 48, 60, 1.0f, 0.0f);
      baseViolaRatio = 1.0f - baseCelloRatio;
    } else if (midiNote <= 72) {
      baseViolaRatio = map(midiNote, 60, 72, 1.0f, 0.0f);
      baseViolinRatio = 1.0f - baseViolaRatio;
    } else {
      baseViolinRatio = 1.0f;
    }
    
    //float bassBlend = baseBassRatio;
    //float celloBlend = baseCelloRatio;
    //float violaBlend = baseViolaRatio;
    //float violinBlend = baseViolinRatio;
    
    float bassBlend = 0.0f;
    float celloBlend = 0.0f;
    float violaBlend = 0.0f;
    float violinBlend = 1.0f;
    
    /*
    // 【特例パッチ】 D6(86) と D#6/Eb6(87) の波形破綻への対応
    if (midiNote == 86 || midiNote == 87) {
      float threshMF = 80.0f / 127.0f;
      float threshFF = 112.0f / 127.0f;
      if (startNorm > threshMF) {
        violinBlend = map(constrain(startNorm, threshMF, threshFF), threshMF, threshFF, 1.0f, 0.0f);
        violaBlend = 1.0f - violinBlend;
      }
      bassBlend = 0.0f; celloBlend = 0.0f;
    }
    */
    
    
    /*
    // ---------------------------------------------------------
    // 2. 新設計：各楽器のプロファイルを「必要な分だけ」動的にオシレーター化して配置
    // ---------------------------------------------------------
    totalActiveOscils = 0; // カウンターリセット
    
    // ループ処理を共通化するための構造化（内部用の一時クラスや処理の連続実行）
    // 比率が 0.0 より大きい楽器だけを狙い撃ちしてオシレーターを生成する（安全装置・負荷対策）
    
    if (bassBlend > 0.0f) {
      float[][] pStart = utils.GetDynamicProfile(midiNote, startVel, "Bass", 24, 0.9f);
      createInstrumentOscillators(pStart, bassBlend, volFactor, globalVibSpeed, fmDepthCents, mixer);
    }
    if (celloBlend > 0.0f) {
      float[][] pStart = utils.GetDynamicProfile(midiNote, startVel, "Cello", 36, 0.9f);
      createInstrumentOscillators(pStart, celloBlend, volFactor, globalVibSpeed, fmDepthCents, mixer);
    }
    if (violaBlend > 0.0f) {
      float[][] pStart = utils.GetDynamicProfile(midiNote, startVel, "Viola", 48, 0.9f);
      createInstrumentOscillators(pStart, violaBlend, volFactor, globalVibSpeed, fmDepthCents, mixer);
    }
    if (violinBlend > 0.0f) {
      float[][] pStart = utils.GetDynamicProfile(midiNote, startVel, "Violin", 55, 0.9f);
      createInstrumentOscillators(pStart, violinBlend, volFactor, globalVibSpeed, fmDepthCents, mixer);
    }
    */
    /* 
    // ---------------------------------------------------------
    // 2. 改修版：各楽器のプロファイルを「動的な上限数」でオシレーター化
    // ---------------------------------------------------------
    // --- コンストラクタ下部：各楽器呼び出しの直前に追加 ---
    //float velRatio = pow(endVel / (startVel + 0.0001f), 0.7f);
    float velRatio = pow((endVel + 0.0001f) / (startVel + 0.0001f), 0.9f);
    
    totalActiveOscils = 0; // カウンターリセット
    
    if (bassBlend > 0.0f) {
      float[][] pStart = utils.GetDynamicProfile(midiNote, startVel, "Bass", 24, 0.9f);
      // ユーザー様提案のルール：1.0未満（ブレンド時）なら20、単一なら30を上限とする
      int maxPeaks = (bassBlend < 1.0f) ? 20 : 30;
      createInstrumentOscillators(pStart, bassBlend, maxPeaks, volFactor, globalVibSpeed, fmDepthCents, mixer, velRatio);
    }
    
    if (celloBlend > 0.0f) {
      float[][] pStart = utils.GetDynamicProfile(midiNote, startVel, "Cello", 36, 0.9f);
      int maxPeaks = (celloBlend < 1.0f) ? 20 : 30;
      createInstrumentOscillators(pStart, celloBlend, maxPeaks, volFactor, globalVibSpeed, fmDepthCents, mixer, velRatio);
    }
    
    if (violaBlend > 0.0f) {
      float[][] pStart = utils.GetDynamicProfile(midiNote, startVel, "Viola", 48, 0.9f);
      int maxPeaks = (violaBlend < 1.0f) ? 20 : 30;
      createInstrumentOscillators(pStart, violaBlend, maxPeaks, volFactor, globalVibSpeed, fmDepthCents, mixer, velRatio);
    }
    
    if (violinBlend > 0.0f) {
      float[][] pStart = utils.GetDynamicProfile(midiNote, startVel, "Violin", 55, 0.9f);
      int maxPeaks = (violinBlend < 1.0f) ? 20 : 30;
      createInstrumentOscillators(pStart, violinBlend, maxPeaks, volFactor, globalVibSpeed, fmDepthCents, mixer, velRatio);
    }
    */
    
    // 更新日6/12
    // =========================================================
    // 【新設計】開始・終了ベロシティから音量変化を正しくスケーリング
    // =========================================================
    // 1. 倍音レシピを安全に取り出すための基準ベロシティ（0除算・無音化の防止）
    float profileVel = (startVel > 0) ? startVel : endVel;
    if (profileVel <= 0) profileVel = 64.0f; // 両方0の場合の完全セーフティ
    
    // 2. 開始時と終了時の実際の音量ゲインファクターを、0.9乗カーブに基づいて個別に計算
    float startGainFactor = pow(startVel / profileVel, 0.9f);
    float endGainFactor   = pow(endVel / profileVel, 0.9f);
    
    // ---------------------------------------------------------
    // 2. 各楽器のプロファイルを動的にオシレーター化して配置
    // ---------------------------------------------------------
    totalActiveOscils = 0; // カウンターリセット
    
    if (bassBlend > 0.0f) {
      float[][] pStart = utils.GetDynamicProfile(midiNote, profileVel, "Bass", 24, 0.9f); // profileVelを渡す
      int maxPeaks = (bassBlend < 1.0f) ? 20 : 30;
      createInstrumentOscillators(pStart, bassBlend, maxPeaks, volFactor, globalVibSpeed, fmDepthCents, mixer, startGainFactor, endGainFactor); // ファクターを渡す
    }
    
    if (celloBlend > 0.0f) {
      float[][] pStart = utils.GetDynamicProfile(midiNote, profileVel, "Cello", 36, 0.9f);
      int maxPeaks = (celloBlend < 1.0f) ? 20 : 30;
      createInstrumentOscillators(pStart, celloBlend, maxPeaks, volFactor, globalVibSpeed, fmDepthCents, mixer, startGainFactor, endGainFactor);
    }
    
    if (violaBlend > 0.0f) {
      float[][] pStart = utils.GetDynamicProfile(midiNote, profileVel, "Viola", 48, 0.9f);
      int maxPeaks = (violaBlend < 1.0f) ? 20 : 30;
      createInstrumentOscillators(pStart, violaBlend, maxPeaks, volFactor, globalVibSpeed, fmDepthCents, mixer, startGainFactor, endGainFactor);
    }
    
    if (violinBlend > 0.0f) {
      float[][] pStart = utils.GetDynamicProfile(midiNote, profileVel, "Violin", 55, 0.9f);
      int maxPeaks = (violinBlend < 1.0f) ? 20 : 30;
      createInstrumentOscillators(pStart, violinBlend, maxPeaks, volFactor, globalVibSpeed, fmDepthCents, mixer, startGainFactor, endGainFactor);
    }
    
    
    /*
    // =========================================================
    // 【データ分析から導いた数式③】 摩擦ノイズ(Bite)の動的パラメータ
    // =========================================================
    // 1. 各楽器の基準値を設定（Python解析データより）
    float biteCutoffBass   = 1150f;  float biteDecayBass   = 0.180f;
    float biteCutoffCello  = 1840f;  float biteDecayCello  = 0.093f;
    float biteCutoffViola  = 3280f;  float biteDecayViola  = 0.116f;
    float biteCutoffViolin = 3370f;  float biteDecayViolin = 0.035f;

    // 2. 現在の音階（クロスフェード割合）に合わせてパラメータをブレンド
    float blendedBiteCutoff = (biteCutoffBass * baseBassRatio) + (biteCutoffCello * baseCelloRatio) 
                            + (biteCutoffViola * baseViolaRatio) + (biteCutoffViolin * baseViolinRatio);
                            
    float blendedBiteDecay  = (biteDecayBass * baseBassRatio) + (biteDecayCello * baseCelloRatio) 
                            + (biteDecayViola * baseViolaRatio) + (biteDecayViolin * baseViolinRatio);

    // グラフの傾向（高音ほどノイズが短く・高くなる）を少し加味する微調整
    // 基音(f0)が上がるにつれて、カットオフを少し上げ、ディケイを少し短くする
    float pitchFactor = map(midiNote, 24, 84, 0.8f, 1.2f);
    float finalBiteCutoff = constrain(blendedBiteCutoff * pitchFactor, 80f, 6000f);
    float finalBiteDecay  = constrain(blendedBiteDecay / pitchFactor, 0.01f, 0.3f);


    // --- 摩擦ノイズ層（Bite）の生成 ---
    // アタックの強さ(ベロシティ)に応じてノイズ音量を変える
    float biteVol = masterVol * 0.001f * startNorm; 
    
    Noise biteNoise = new Noise(biteVol, Noise.Tint.WHITE);
    MoogFilter biteFilter = new MoogFilter(finalBiteCutoff, 0.2f); // 動的カットオフを適用
    biteFilter.type = MoogFilter.Type.LP; 
    
    // Attackは固定(超短く)、Decayに動的パラメータを適用
    biteEnv = new ADSR(1.0f, 0.1f, finalBiteDecay, 0.0f, 0.1f);
    
    biteNoise.patch(biteFilter);
    biteFilter.patch(biteEnv);
    // (biteEnvは直接outへ出すためmixerには繋がない)
    */
    
    
    /* 【持続ノイズHiss削除　6/12】
    // --- 共鳴ノイズ層（Hiss） ---
    startHissVol = masterVol * 0.001f * startNorm;
    endHissVol   = masterVol * 0.001f * endNorm;
    
    Noise hissNoise = new Noise(0.0f, Noise.Tint.PINK);
    hissAmpLine = new Line();
    hissAmpLine.patch(hissNoise.amplitude);
    
    Oscil hissWobble = new Oscil(globalVibSpeed, startHissVol * 0.3f, Waves.SINE);
    hissWobble.setPhase(random(1.0f));
    hissWobble.patch(hissNoise.amplitude);
    
    MoogFilter hissFilter = new MoogFilter(1500f, 0.15f);
    hissFilter.type = MoogFilter.Type.HP;
    hissEnv = new ADSR(1.0f, 0.05f, 0.0f, 1.0f, 0.1f);
    
    hissNoise.patch(hissFilter);
    hissFilter.patch(hissEnv); 
    hissEnv.patch(mixer);
    */
  }
  
  /*
  // 各楽器のピークデータをグローバルなオシレーター配列に安全に展開するヘルパー関数
  private void createInstrumentOscillators(float[][] profile, float blendRatio, float volFactor, float globalVibSpeed, float fmDepthCents, Summer mixer) {
    if (profile == null) return;
    
    byte numPeaks = (byte)profile.length;
    for (byte h = 0; h < numPeaks; h++) {
      // 安全装置：配列の上限を超えないようにガード
      if (totalActiveOscils >= MAX_TOTAL_OSCILS) break;
      
      float targetFreq = profile[h][0]; // Pythonが実測した、非整数倍音が含まれる生の周波数
      float rawAmp     = profile[h][1]; // 正規化された振幅
      
      // 20000Hz超えの超音波は生成をスキップしてCPUを徹底節約（アンチエイリアス）
      if (targetFreq > 22050.0f) continue;
      
      // 【核心】元のスペクトル振幅に、全体のボリュームと「この楽器のブレンド比率」を掛け合わせる
      float finalStartAmp = rawAmp * volFactor * blendRatio;
      float finalEndAmp   = finalStartAmp * (endHissVol / (startHissVol + 0.0001f)); // 代案2の音量比率の維持
      
      int idx = totalActiveOscils;
      startAmps[idx] = finalStartAmp;
      endAmps[idx]   = finalEndAmp;
      
      // オシレーターとエンベロープの生成
      oscils[idx] = new Oscil(targetFreq, 0.0f, Waves.SINE);
      oscils[idx].setPhase(0.0f);
      
      // ビブラート（FM変調）の適用
      float vibDepthRatio = pow(2.0f, fmDepthCents / 1200.0f) - 1.0f;
      targetVibDepths[idx] = targetFreq * vibDepthRatio;
      
      Oscil freqController = new Oscil(globalVibSpeed, 0.0f, Waves.SINE);
      freqController.offset.setLastValue(targetFreq);
      freqController.patch(oscils[idx].frequency);
      
      vibDepthLines[idx] = new Line();
      vibDepthLines[idx].patch(freqController.amplitude);
      
      // 音量揺らぎ（AM変調）の適用 (高次倍音ほど激しく揺らす特性は継承)
      float wobbleDepth = map(idx, 0, MAX_TOTAL_OSCILS - 1, 0.15f, 0.60f);
      float wobbleAmp = startAmps[idx] * wobbleDepth;
      
      Oscil ampLFO = new Oscil(globalVibSpeed, wobbleAmp, Waves.SINE);
      ampLFO.setPhase(random(1.0f));
      ampLFO.patch(oscils[idx].amplitude);
      
      ampLines[idx] = new Line();
      ampLines[idx].patch(oscils[idx].amplitude);
      
      oscils[idx].patch(mixer);
      
      totalActiveOscils++; // 生成に成功したら総数をインクリメント
    }
  }
  */
  /*
  // 各楽器のピークデータを、振幅選別フィルタを通した上で安全にオシレーター展開するヘルパー関数
  private void createInstrumentOscillators(final float[][] profile, float blendRatio, int maxPeaks, float volFactor, float globalVibSpeed, float fmDepthCents, Summer mixer, float velRatio) {
    if (profile == null || profile.length == 0) return;
    
    int rawLength = profile.length;
    // 実際に採用する山の数を決定（データ数と制限数の小さい方）
    final int limitPeaks = min(maxPeaks, rawLength);
    
    // ---------------------------------------------------------
    // 【振幅選別フィルタ】振幅が大きい順にインデックスをソートする
    // ---------------------------------------------------------
    Integer[] indices = new Integer[rawLength];
    for (int i = 0; i < rawLength; i++) indices[i] = i;
    
    Arrays.sort(indices, new Comparator<Integer>() {
      public int compare(Integer a, Integer b) {
        // 振幅[1]の降順（大きい順）で並び替え
        return Float.compare(profile[b][1], profile[a][1]);
      }
    });
    
    // 上位 N 個（limitPeaks分）だけを一度抽出し、一時配列に格納
    float[][] filteredProfile = new float[limitPeaks][2];
    for (int i = 0; i < limitPeaks; i++) {
      filteredProfile[i][0] = profile[indices[i]][0]; // 周波数
      filteredProfile[i][1] = profile[indices[i]][1]; // 振幅
    }
    
    // ---------------------------------------------------------
    // 【周波数再整列】Processing側（ビブラート計算など）の仕様に合わせ、周波数の低い順に戻す
    // ---------------------------------------------------------
    Arrays.sort(filteredProfile, new Comparator<float[]>() {
      public int compare(float[] a, float[] b) {
        // 周波数[0]の昇順（低い順）で並び替え
        return Float.compare(a[0], b[0]);
      }
    });
    
    // ---------------------------------------------------------
    // 3. 厳選されたピークデータのみをオシレーターとして生成・接続
    // ---------------------------------------------------------
    for (int h = 0; h < limitPeaks; h++) {
      if (totalActiveOscils >= MAX_TOTAL_OSCILS) break;
      
      float targetFreq = filteredProfile[h][0]; // 周波数
      float rawAmp     = filteredProfile[h][1]; // 厳選された強い振幅
      
      if (targetFreq > 22050.0f) continue;
      
      float finalStartAmp = rawAmp * volFactor * blendRatio;
      float finalEndAmp   = finalStartAmp * velRatio;
      
      int idx = totalActiveOscils;
      startAmps[idx] = finalStartAmp;
      endAmps[idx]   = finalEndAmp;
      
      // オシレーターとビブラート（LFO）、AM変調の適用（既存の高品質ロジックをそのまま継承）
      oscils[idx] = new Oscil(targetFreq, 0.0f, Waves.SINE);
      oscils[idx].setPhase(0.0f);
      
      float vibDepthRatio = pow(2.0f, fmDepthCents / 1200.0f) - 1.0f;
      targetVibDepths[idx] = targetFreq * vibDepthRatio;
      
      Oscil freqController = new Oscil(globalVibSpeed, 0.0f, Waves.SINE);
      freqController.offset.setLastValue(targetFreq);
      freqController.patch(oscils[idx].frequency);
      
      vibDepthLines[idx] = new Line();
      vibDepthLines[idx].patch(freqController.amplitude);
      
      float wobbleDepth = map(idx, 0, MAX_TOTAL_OSCILS - 1, 0.15f, 0.60f);
      float wobbleAmp = startAmps[idx] * wobbleDepth;
      
      Oscil ampLFO = new Oscil(globalVibSpeed, wobbleAmp, Waves.SINE);
      ampLFO.setPhase(random(1.0f));
      ampLFO.patch(oscils[idx].amplitude);
      
      ampLines[idx] = new Line();
      ampLines[idx].patch(oscils[idx].amplitude);
      
      oscils[idx].patch(mixer);
      
      totalActiveOscils++;
    }
  }
  */
  // 更新日6/12
  private void createInstrumentOscillators(final float[][] profile, float blendRatio, int maxPeaks, float volFactor, float globalVibSpeed, float fmDepthCents, Summer mixer, float startGainFactor, float endGainFactor) {
    if (profile == null || profile.length == 0) return;
    
    int rawLength = profile.length;
    final int limitPeaks = min(maxPeaks, rawLength);
    
    Integer[] indices = new Integer[rawLength];
    for (int i = 0; i < rawLength; i++) indices[i] = i;
    
    Arrays.sort(indices, new Comparator<Integer>() {
      public int compare(Integer a, Integer b) {
        return Float.compare(profile[b][1], profile[a][1]);
      }
    });
    
    float[][] filteredProfile = new float[limitPeaks][2];
    for (int i = 0; i < limitPeaks; i++) {
      filteredProfile[i][0] = profile[indices[i]][0]; 
      filteredProfile[i][1] = profile[indices[i]][1]; 
    }
    
    Arrays.sort(filteredProfile, new Comparator<float[]>() {
      public int compare(float[] a, float[] b) {
        return Float.compare(a[0], b[0]);
      }
    });
    
    for (int h = 0; h < limitPeaks; h++) {
      if (totalActiveOscils >= maxTotalOscils) break;
      
      float targetFreq = filteredProfile[h][0]; 
      float rawAmp     = filteredProfile[h][1]; 
      
      //if (targetFreq > 22050.0f) continue;
      if (targetFreq > 14000.0f) continue;
      
      // 6/13【新設：数学的倍音フィルター】
      // 8000Hzから18000Hzにかけて、その倍音の振幅（パワー）を滑らかに 100% から 0% へ減衰させる
      if (targetFreq > 7000.0f) {
        float filterFactor = map(targetFreq, 7000.0f, 17000.0f, 1.0f, 0.0f);
        rawAmp *= constrain(filterFactor, 0.0f, 1.0f); // 厳密に0〜1の範囲に固定して乗算
      }
      
      // 【確実なスケーリング】実際の開始音量と終了音量を個別に決定
      float baseAmp = rawAmp * volFactor * blendRatio;
      float finalStartAmp = baseAmp * startGainFactor;
      float finalEndAmp   = baseAmp * endGainFactor;
      
      int idx = totalActiveOscils;
      startAmps[idx] = finalStartAmp;
      endAmps[idx]   = finalEndAmp;
      
      oscils[idx] = new Oscil(targetFreq, 0.0f, Waves.SINE);
      oscils[idx].setPhase(random(1.0f));
      
      // --- ビブラート（FM変調）の適用 ---
      float vibDepthRatio = pow(2.0f, fmDepthCents / 1200.0f) - 1.0f;
      targetVibDepths[idx] = targetFreq * vibDepthRatio;
      
      Oscil freqController = new Oscil(globalVibSpeed, 0.0f, Waves.SINE);
      freqController.offset.setLastValue(targetFreq);
      //freqController.patch(freqController.amplitude); // ※元のコードママ
      freqController.patch(oscils[idx].frequency);
      
      vibDepthLines[idx] = new Line();
      vibDepthLines[idx].patch(freqController.amplitude);
      
      // --- 音量揺らぎ（AM変調）：ウナリ自体をフェードさせる革新的な設計 ---
      wobbleLines[idx] = new Line();
      Oscil ampLFO = new Oscil(globalVibSpeed, 0.0f, Waves.SINE);
      ampLFO.setPhase(random(1.0f));
      
      wobbleLines[idx].patch(ampLFO.amplitude); // 新設LineをLFOの振幅へパッチ！これでウナリがフェードする
      ampLFO.patch(oscils[idx].amplitude);      // LFOの出力をオシレーターへ（加算）
      
      // ベース音量の適用
      ampLines[idx] = new Line();
      ampLines[idx].patch(oscils[idx].amplitude); // ベース音量変化をオシレーターへ（加算）
      
      oscils[idx].patch(mixer);
      totalActiveOscils++;
    }
  }
  
  
  
  /*
  void noteOn(float dur) {
    masterEnv.noteOn();
    biteEnv.noteOn(); 
    hissEnv.noteOn();
    
    float safeDur = dur + 0.7f; 
    // ビブラートが最大になるまでの時間（0.5秒、または音符の長さの短い方）
    float vibDelayTime = min(0.5f, dur); 
    
    
    //for(int i = 0; i < numHarmonics; i++) {
      // 【修正】インスタンスが生成されていない（null）場合は、それ以上の倍音はないのでループを終了する
    //  if (ampLines[i] == null || vibDepthLines[i] == null) {
    //    continue; 
    //  }
    //  
    //  
    //  ampLines[i].activate(safeDur, startAmps[i], endAmps[i]);
    //  // 【実行】0Hzから目標の揺れ幅(Hz)へ、vibDelayTimeかけて徐々にビブラートを深くする
    //  vibDepthLines[i].activate(vibDelayTime, 0.0f, targetVibDepths[i]);
    //}
    
    // noteOn 内のループ部分の修正
    for(int i = 0; i < totalActiveOscils; i++) {
      if (ampLines[i] == null || vibDepthLines[i] == null) {
        continue; 
      }
      ampLines[i].activate(safeDur, startAmps[i], endAmps[i]);
      vibDepthLines[i].activate(vibDelayTime, 0.0f, targetVibDepths[i]);
    }
    
    hissAmpLine.activate(safeDur, startHissVol, endHissVol);
    
    masterEnv.patch(out); 
    biteEnv.patch(out); 
  }
  */
  
  // 更新日6/12
  void noteOn(float dur) {
    masterEnv.noteOn();
    //biteEnv.noteOn(); 
    //hissEnv.noteOn();
    
    float safeDur = dur + 0.7f; 
    float vibDelayTime = min(0.5f, dur); 
    
    for(int i = 0; i < totalActiveOscils; i++) {
      if (ampLines[i] == null || vibDepthLines[i] == null || wobbleLines[i] == null) {
        continue; 
      }
      // 1. ベース音量のフェード
      ampLines[i].activate(safeDur, startAmps[i], endAmps[i]);
      
      // 2. ウナリ（AM変調）の揺れ幅フェード（ベース音量に比例してウナリの深さも綺麗に消える）
      float wobbleDepth = map(i, 0, maxTotalOscils - 1, 0.15f, 0.60f);
      wobbleLines[i].activate(safeDur, startAmps[i] * wobbleDepth, endAmps[i] * wobbleDepth);
      
      // 3. ビブラート（FM変調）の変化
      vibDepthLines[i].activate(vibDelayTime, 0.0f, targetVibDepths[i]);
    }
    
    //hissAmpLine.activate(safeDur, startHissVol, endHissVol);
    
    masterEnv.patch(longOut); 
    //biteEnv.patch(longOut); 
  }
  
  void noteOff() {
    masterEnv.noteOff();
    //for(int i=0; i<numHarmonics; i++) harmonicEnvs[i].noteOff();
    //biteEnv.noteOff(); 
    //hissEnv.noteOff();
    masterEnv.unpatchAfterRelease(longOut); 
    //biteEnv.unpatchAfterRelease(longOut);
  }
}
\end{lstlisting}

\subsection{Piano.pde}
\begin{lstlisting}[caption={Piano.pde}, label={lst:proc_piano}]

//float[][] pianoDecayLUT;
//float[][] pianoWobbleRateLUT;
//float[][] pianoWobbleDepthLUT;

// ==========================================
// 究極のピアノ・インストゥルメント (デュアル・ストリング & ハイブリッドデータ駆動版)
// ==========================================
class PianoInstrument implements Instrument {
  // 分析データに合わせて最大80倍音まで生成
  byte numHarmonics/* = 40*/;
  
  ADSR masterEnv; // 全体の音の長さを管理し、離鍵時にミュートする（ダンパーフェルトの役割）
  ADSR hammerEnv; // ハンマーが弦を叩く一瞬のノイズ用エンベロープ
  
  // 個別の減衰（Damp）を管理するリスト。
  // 芯の弦と共鳴弦でDampの数が変動するため、配列ではなくArrayListを使用します。
  ArrayList<Damp> damps = new ArrayList<Damp>();
  
  PianoInstrument(float midiNote, float velocity, float masterVol) {
    Summer mixer = new Summer();
    
    // ユーザー定義のベロシティ閾値
    float thresholdPP = 32.0f / 127.0f;
    float thresholdMF = 80.0f / 127.0f;
    float thresholdFF = 112.0f / 127.0f;
    
    // 現在のベロシティ層に応じて、ベースとなる波形プロファイル（周波数・振幅の骨組み）を決定
    String baseVelocity;
    float t = 0;
    //boolean needInterpolation = false;
    //float[][] lowerProfile = null;
    //float[][] upperProfile = null;
    
    if (velocity <= thresholdPP) {
      baseVelocity = "PP";
    } else if (velocity <= thresholdMF) {
      // pp ～ mf の間：構造が近い方のスペクトル形状を選択（あるいは、より高度に処理するためのフラグ設定）
      t = map(velocity, thresholdPP, thresholdMF, 0.0f, 1.0f);
      baseVelocity = (t < 0.5f) ? "PP" : "MF";
    } else if (velocity <= thresholdFF) {
      // mf ～ ff の間
      t = map(velocity, thresholdMF, thresholdFF, 0.0f, 1.0f);
      baseVelocity = (t < 0.5f) ? "MF" : "FF";
    } else {
      baseVelocity = "FF";
    }
    
    float[][]pianoDecayLUT = utils.GetLUT("PianoDecay" + baseVelocity);
    float[][]pianoWobbleRateLUT = utils.GetLUT("PianoWobbleRate" + baseVelocity);
    float[][]pianoWobbleDepthLUT = utils.GetLUT("PianoWobbleDepth" + baseVelocity);
    
    // ==========================================
    // 1. 全体のエンベロープ設定（ダンパーの挙動）
    // ==========================================
    // 鍵盤を押している間は減衰を邪魔せず(Sustain=1.0)、離した瞬間(noteOff)に0.15秒でミュート
    masterEnv = new ADSR(1.0f, 0.001f, 0.4f, 0.5f, 0.09f);
    mixer.patch(masterEnv); 
    
    //float f0 = Frequency.ofMidiNote(midiNote).asHz();
    float normVel = pow(constrain(velocity, 0, 127) / 127.0f, 0.9f);
    float volFactor = masterVol * 0.022f; 
    
    //volFactor *= (0.5 + (midiNote/254.0f));  // 必要かどうか検討 --> 不要
    
    // ==========================================
    // 2. LUTデータの取得とインハーモニシティ設定
    // ==========================================
    // 強弱に応じた倍音構成（音色レシピ）を取得
    float[][] profile = utils.GetDynamicProfile(midiNote, velocity, "Piano", 21, 0.9f);
    
    // 【修正の核心1】実測ピークデータから、今回のループ回数を動的に確定させる！
    numHarmonics = (byte)profile.length;
    //if (numHarmonics > 30) numHarmonics = 30;
    
    /*
    // インハーモニシティ（弦の硬さによるピッチの上擦り）。高音ほどズレが大きくなる。
    float B = 0.0001f * pow(f0 / 261.6f, 1.5f); 
    B = constrain(B, 0.00002f, 0.003f);     
    */
    
    // 現在の音階のLUT配列用インデックス (A0(MIDI:21)を0とする)
    int lutIndex = constrain(round(midiNote) - 21, 0, 87);
    
    // Pythonで解析したJSONデータが欠損していないかチェック
    // ※汎用メソッド LoadLUT で読み込んだ3つのグローバル配列を使用します
    //boolean hasRawData = (pianoDecayLUT != null && pianoDecayLUT[lutIndex] != null);
    // 【修正の核心2】コンパイル警告を完全に消し去るための厳格な3軸構造チェック
    /*boolean hasRawData = false;
    if (pianoDecayLUT != null && lutIndex < pianoDecayLUT.length && pianoDecayLUT[lutIndex] != null) {
      if (pianoWobbleRateLUT != null && lutIndex < pianoWobbleRateLUT.length && pianoWobbleRateLUT[lutIndex] != null) {
        if (pianoWobbleDepthLUT != null && lutIndex < pianoWobbleDepthLUT.length && pianoWobbleDepthLUT[lutIndex] != null) {
          hasRawData = true;
        }
      }
    }*/

    // --- データ欠損時のバックアップ用（数式モデル）の事前計算 ---
    // Pythonのグラフ解析から導き出した「ピアノの物理特性の傾向」を数式化
    float mathBaseDecay = 29.0f * pow(0.9757f, midiNote);
    if (midiNote < 40) mathBaseDecay = min(mathBaseDecay, 20.0f);            // 低音の物理的限界（頭打ち）
    if (midiNote >= 89) mathBaseDecay *= map(midiNote, 89, 108, 1.0f, 3.0f); // 超高音のダンパーレス構造による伸び
    
    float mathBaseWobbleRate = map(abs(midiNote - 66.0f), 0, 45, 0.5f, 4.5f); // うなりの速さ(MIDI66を底としたU字型)
    float mathBaseDepth = map(midiNote, 21, 108, 0.03f, 0.15f);               // うなりの深さ(高音ほど深い)

    // リバーブ（共鳴）が最大音量になるまでの時間。低音はゆっくり、高音は速く立ち上がる。
    float buildUpTime = map(midiNote, 21, 108, 3.0f, 0.5f);

    // ==========================================
    // 3. 弦のモデリング（デュアル・ストリング構造）
    // ==========================================
    for (int i = 0; i < numHarmonics; i++) {
      float amp = profile[i][1] * volFactor;
      if (amp <= 0.0001f) continue; // 音量が極端に小さい倍音はCPU負荷軽減のためスキップ
      // 省エネ設計
      
      /*
      int h = i + 1;
      float inharmonicity = sqrt(1.0f + B * h * h);
      float targetFreq = f0 * h * inharmonicity; // 実際の周波数
      */
      
      float targetFreq = profile[i][0]; // 実測された上擦り済みの正確な周波数
      if (targetFreq > 22050.0f) continue; // 人間の可聴域を超えたらスキップ
      
      float decayTime, wobbleRate, wobbleDepth;
      
      // 実測データ(JSON)があれば使い、なければ数式で補完する「ハイブリッド方式」
      //if (hasRawData && i < pianoDecayLUT[lutIndex].length && pianoDecayLUT[lutIndex][i] > 0.1f) {
      // 実測データ(JSON)の適用と数式補完
      if (pianoDecayLUT != null && pianoWobbleRateLUT != null && pianoWobbleDepthLUT != null && 
          pianoDecayLUT[lutIndex] != null &&
          i < pianoDecayLUT[lutIndex].length) {
        decayTime = pianoDecayLUT[lutIndex][i];
        wobbleRate = pianoWobbleRateLUT[lutIndex][i];
        wobbleDepth = pianoWobbleDepthLUT[lutIndex][i];
        
        // もしデータが0、または異常値だった場合は数式モデルで安全に保護する
        if (decayTime <= 0.1f) {
          float harmonicDecayCurve = (i <= 14) ? map(i, 0, 14, 1.0f, 0.3f) : map(i, 14, 80, 0.3f, 1.5f);
          decayTime = max(mathBaseDecay * harmonicDecayCurve, 0.1f);
        }
      } else {
        // データがない場合は、インデックス14を底とする減衰のU字カーブを適用
        float harmonicDecayCurve = (i <= 14) ? map(i, 0, 14, 1.0f, 0.3f) : map(i, 14, 80, 0.3f, 1.5f);
        decayTime = max(mathBaseDecay * harmonicDecayCurve, 0.1f);
        wobbleRate = mathBaseWobbleRate * (1.0f + i * 0.02f); 
        wobbleDepth = mathBaseDepth * pow(0.9f, i);
      }
      
      decayTime *= 0.7;
      
      // ----------------------------------------
      // レイヤーA：アタックの芯（Dry弦）
      // ----------------------------------------
      Oscil coreOsc = new Oscil(targetFreq, amp, Waves.SINE);
      // アタックのインパクトを最大化するため、位相(波のスタート位置)は0付近に固定
      coreOsc.setPhase(random(0.0f, 0.1f)); 
      
      // 0.001秒で素早く立ち上がり、それぞれの寿命(decayTime)で消えていく
      Damp coreDamp = new Damp(0.001f, decayTime);
      coreOsc.patch(coreDamp).patch(mixer);
      damps.add(coreDamp);
      
      // ----------------------------------------
      // レイヤーB：時間差で来る響き（共鳴弦）
      // ----------------------------------------
      float reverbAmp = amp * wobbleDepth; // 共鳴の音量は、芯の音の数%〜十数%程度に抑える
      
      // うなりのスピード(wobbleRate)が設定されており、かつ音量が十分ある場合のみ共鳴弦を生成
      if (reverbAmp > 0.0001f && wobbleRate > 0.0f) {
        // 周波数をわずかにズラす（Detune）ことで、芯の弦と干渉して自然な「うなり(Beat)」を生む
        Oscil resOsc = new Oscil(targetFreq + wobbleRate, reverbAmp, Waves.SINE);
        //resOsc.setPhase(random(1.0f)); // 芯の弦と複雑に干渉させるため、位相は完全にランダム
        
        // アタックタイムを buildUpTime に設定し、数秒かけてゆっくり音量を上げる
        Damp resDamp = new Damp(buildUpTime, decayTime);
        resOsc.patch(resDamp).patch(mixer);
        damps.add(resDamp);
      }
    }
    
    // ==========================================
    // 4. ハンマーの打撃ノイズ（アタック）
    // ==========================================
    float hammerVol = masterVol * 0.003f * normVel; 
    Noise hammerNoise = new Noise(hammerVol, Noise.Tint.WHITE);
    
    // 高音の鍵盤ほど、硬い「カチッ」としたノイズになるようにフィルターを調整
    float hammerCutoff = map(midiNote, 21, 108, 600f, 4000f);
    MoogFilter hammerFilter = new MoogFilter(hammerCutoff, 0.1f);
    hammerFilter.type = MoogFilter.Type.LP; 
    
    // 0.03秒で消滅する、極めてタイトなエンベロープ（ミュートの影響を受けない）
    hammerEnv = new ADSR(1.0f, 0.001f, 0.02f, 0.0f, 0.01f);
    
    hammerNoise.patch(hammerFilter).patch(hammerEnv);
  }
  
  // ==========================================
  // 発音と停止のコントロール
  // ==========================================
  void noteOn(float dur) {
    masterEnv.noteOn();
    hammerEnv.noteOn();
    
    // 準備した芯の弦と共鳴弦のDamp（エンベロープ）を一斉にスタート
    for (Damp d : damps) {
      d.activate(); 
    }
    
    masterEnv.patch(longOut); 
    hammerEnv.patch(longOut); // ハンマーノイズは独立して出力
  }
  
  void noteOff() {
    masterEnv.noteOff(); // ここでダンパーペダルが降りてミュート開始
    hammerEnv.noteOff(); 
    masterEnv.unpatchAfterRelease(longOut); 
    hammerEnv.unpatchAfterRelease(longOut);
  }
}
\end{lstlisting}

\subsection{Drum.pde}
\begin{lstlisting}[caption={Drum.pde}, label={lst:proc_drum}]
// ドラム用LUTの宣言 [インデックス][0:周波数, 1:強度, 2:衰退時間(Decay)]
//float[][] kickLUT, snareLUT, hihatLUT;











// ==========================================
// 1. キック（Kick）のクラス
// ==========================================
class KickInstrument implements Instrument {
  //Summer mixer;
  ADSR masterEnv;
  
  //Oscil[] tones;
  ADSR[] harmonicEnvs;
  //Line[] pitchSweeps;

  KickInstrument(float masterVol, float velocity) {
    Summer mixer = new Summer();
    
    float[][]kickLUT = utils.GetLUT("Kick");
    
    // フルートと同じく、リリース時のバッファを持たせたマスターエンベロープ
    masterEnv = new ADSR(1.0f, 0.001f, 0.0f, 1.0f, 0.1f);
    mixer.patch(masterEnv);
    
    //int numComponents = kickLUT.length;
    int numComponents = 50;
    Oscil[] tones = new Oscil[numComponents];
    harmonicEnvs = new ADSR[numComponents];
    Line[] pitchSweeps = new Line[numComponents];
    
    //float baseAmp = masterVol * pow(constrain(velocity, 0, 127) / 127.0f, 0.9f) * 25.0f;
    float baseAmp = masterVol * pow(constrain(velocity, 0, 127) / 127.0f, 0.9f) * 14.0f;

    for (int i = 0; i < numComponents; i++) {
      if (kickLUT[i] == null || kickLUT[i].length < 3) continue;
      
      float freq = kickLUT[i][0] * random(0.95f, 1.05f);
      float mag = kickLUT[i][1];
      float decayTime = kickLUT[i][2];
      float amp = (baseAmp * mag) / numComponents;
      
      tones[i] = new Oscil(freq, amp, Waves.SINE);
      //tones[i].setPhase(random(1.0f)); // 初期位相を散らしてアタックのピーク割れを防ぐ
      
      // 個別の成分のエンベロープ
      //harmonicEnvs[i] = new ADSR(1.0f, 0.001f, decayTime * 0.05f, 0.0f, 0.1f);
      harmonicEnvs[i] = new ADSR(1.0f, 0.001f, decayTime*0.5f, 0.0f, 0.1f);
      
      // ピッチスイープ
      pitchSweeps[i] = new Line(0.01f, freq * 1.5f, freq);
      pitchSweeps[i].patch(tones[i].frequency);
      
      // 結線: Oscil -> 個別ADSR -> Mixer
      tones[i].patch(harmonicEnvs[i]);
      harmonicEnvs[i].patch(mixer);
    }
  }

  void noteOn(float duration) {
    masterEnv.noteOn();
    for (int i = 0; i < harmonicEnvs.length; i++) {
      if (harmonicEnvs[i] != null) harmonicEnvs[i].noteOn();
    }
    // グローバルな out に直接パッチ！
    masterEnv.patch(shortOut);
  }

  void noteOff() {
    masterEnv.noteOff();
    for (int i = 0; i < harmonicEnvs.length; i++) {
      if (harmonicEnvs[i] != null) harmonicEnvs[i].noteOff();
    }
    // グローバルな out からアンパッチ！
    masterEnv.unpatchAfterRelease(shortOut);
  }
}

// ==========================================
// 2. スネア（Snare）のクラス
// ==========================================
class SnareInstrumentOld implements Instrument {
  //Summer mixer;
  ADSR masterEnv;
  
  // LUT（胴鳴り・倍音）用
  //Oscil[] tones;
  ADSR[] harmonicEnvs;
  //Line[] pitchSweeps;
  
  // ノイズ（スナッピー）用
  //Noise snareNoise;
  //HighPassSP snareFilter;
  ADSR noiseEnv;

  SnareInstrumentOld(float masterVol, float velocity) {
    Summer mixer = new Summer();
    masterEnv = new ADSR(1.0f, 0.001f, 0.0f, 1.0f, 0.1f);
    mixer.patch(masterEnv);
    
    float baseAmp = masterVol * pow(constrain(velocity, 0, 127) / 127.0f, 0.9f) * 4.0f;
    
    float[][]snareLUT = utils.GetLUT("Snare");
    
    // --- 1. 胴鳴りパート (LUTによる加算合成) ---
    int numComponents = snareLUT.length;
    Oscil[] tones = new Oscil[numComponents];
    harmonicEnvs = new ADSR[numComponents];
    Line[] pitchSweeps = new Line[numComponents];

    for (int i = 0; i < numComponents; i++) {
      if (snareLUT[i] == null || snareLUT[i].length < 3) continue;
      
      float freq = snareLUT[i][0];
      float mag = snareLUT[i][1];
      float decayTime = snareLUT[i][2];
      float amp = (baseAmp * mag) / numComponents;
      
      tones[i] = new Oscil(freq, amp, Waves.SINE);
      tones[i].setPhase(random(1.0f)); 
      
      harmonicEnvs[i] = new ADSR(1.0f, 0.001f, decayTime * 0.7, 0.0f, 0.1f);
      
      // スネア特有：低音（胴鳴り）のみピッチスイープ
      if (freq < 1000) {
        pitchSweeps[i] = new Line(0.05f, freq * 1.5f, freq);
        pitchSweeps[i].patch(tones[i].frequency);
      }
      
      tones[i].patch(harmonicEnvs[i]);
      harmonicEnvs[i].patch(mixer);
    }
    
    // --- 2. スナッピーパート (ノイズによる減算合成) ---
    // 以前の解析で導き出した数値をそのまま採用
    float noiseVol = baseAmp * 0.022f; 
    Noise snareNoise = new Noise(noiseVol, Noise.Tint.WHITE);
    HighPassSP snareFilter = new HighPassSP(512, out.sampleRate()); // 7750Hz以上を通す
    noiseEnv = new ADSR(1.0f, 0.003f, 0.15f, 0.0f, 0.1f); // ディケイ0.25秒
    
    // ノイズをフィルターし、エンベロープを通してミキサーへ合流
    snareNoise.patch(snareFilter).patch(noiseEnv).patch(mixer);
  }

  void noteOn(float duration) {
    masterEnv.noteOn();
    for (int i = 0; i < harmonicEnvs.length; i++) {
      if (harmonicEnvs[i] != null) harmonicEnvs[i].noteOn();
    }
    noiseEnv.noteOn();
    masterEnv.patch(shortOut);
  }

  void noteOff() {
    masterEnv.noteOff();
    for (int i = 0; i < harmonicEnvs.length; i++) {
      if (harmonicEnvs[i] != null) harmonicEnvs[i].noteOff();
    }
    noiseEnv.noteOff();
    masterEnv.unpatchAfterRelease(shortOut);
  }
}

// ==========================================
// 3. クローズハイハット（Closed Hi-Hat）のクラス
// ==========================================
class HiHatInstrumentOld implements Instrument {
  //Summer mixer;
  ADSR masterEnv;
  
  //Oscil[] tones;
  ADSR[] harmonicEnvs;
  
  //Noise hatNoise;
  //HighPassSP hatFilter;
  ADSR noiseEnv;

  HiHatInstrumentOld(float masterVol, float velocity) {
    Summer mixer = new Summer();
    // ハイハットは余韻が非常に短いため、マスターのリリースも極短に
    masterEnv = new ADSR(1.0f, 0.001f, 0.0f, 1.0f, 0.05f);
    mixer.patch(masterEnv);
    
    float baseAmp = masterVol * pow(constrain(velocity, 0, 127) / 127.0f, 0.9f) * 5.0f;
    
    float[][]hihatLUT = utils.GetLUT("HiHat");
    
    // --- 1. 金属共鳴パート (LUTによる加算合成) ---
    int numComponents = hihatLUT.length;
    Oscil[] tones = new Oscil[numComponents];
    harmonicEnvs = new ADSR[numComponents];

    for (int i = 0; i < numComponents; i++) {
      if (hihatLUT[i] == null || hihatLUT[i].length < 3) continue;
      
      float freq = hihatLUT[i][0];
      float mag = hihatLUT[i][1];
      float decayTime = hihatLUT[i][2];
      float amp = (baseAmp * mag) / numComponents;
      
      tones[i] = new Oscil(freq, amp, Waves.SINE);
      tones[i].setPhase(random(1.0f));
      
      harmonicEnvs[i] = new ADSR(1.0f, 0.001f, decayTime, 0.0f, 0.05f);
      
      tones[i].patch(harmonicEnvs[i]);
      harmonicEnvs[i].patch(mixer);
    }
    
    // --- 2. 打撃摩擦パート (ノイズによる減算合成) ---
    float noiseVol = baseAmp * 0.042f;
    Noise hatNoise = new Noise(noiseVol, Noise.Tint.WHITE);
    HighPassSP hatFilter = new HighPassSP(10250, out.sampleRate()); // 10000Hz以上の超高音のみ
    noiseEnv = new ADSR(1.0f, 0.001f, 0.06f, 0.0f, 0.05f); // 0.06秒でスパッと消える
    
    hatNoise.patch(hatFilter).patch(noiseEnv).patch(mixer);
  }

  void noteOn(float duration) {
    masterEnv.noteOn();
    for (int i = 0; i < harmonicEnvs.length; i++) {
      if (harmonicEnvs[i] != null) harmonicEnvs[i].noteOn();
    }
    noiseEnv.noteOn();
    masterEnv.patch(shortOut);
  }

  void noteOff() {
    masterEnv.noteOff();
    for (int i = 0; i < harmonicEnvs.length; i++) {
      if (harmonicEnvs[i] != null) harmonicEnvs[i].noteOff();
    }
    noiseEnv.noteOff();
    masterEnv.unpatchAfterRelease(shortOut);
  }
}









// ==========================================
// 2. 新・スネア（Snare）のクラス (シンセサイズ方式)
// ==========================================
class SnareInstrument implements Instrument {
  ADSR masterEnv;
  
  // 1. 胴鳴り（Body）パート用
  Oscil bodyOsc;
  Line  bodyPitchEnv;
  ADSR  bodyAmpEnv;
  
  // 2. 倍音（Overtone）パート用
  Oscil overtoneOsc;
  ADSR  overtoneAmpEnv;
  
  // 3. スナッピー（Noise）パート用
  Noise      snappyNoise;
  //HighPassSP snappyFilter;
  ADSR       snappyAmpEnv;

  SnareInstrument(float masterVol, float velocity, float sinAmp) {
    Summer mixer = new Summer();
    
    // 全体の音量調整
    float baseAmp = masterVol * pow(constrain(velocity, 0, 127) / 127.0f, 0.9f) * 0.15f; 
    float sinSet = constrain(sinAmp, 0, 127) / 127.0f;
    
    // ==========================================
    // パラメータ設定（ここをいじって音色をチューニングします）
    // ==========================================
    
    // --- 1. 胴鳴り（ベースとなる太さ） ---
    float bodyFreq     = 172.3f * random(0.99f, 1.01f);  // ★データ最上部の「172.3 Hz (強さ: 0.63)」をそのまま指定！
    float bodyPitchMod = 3.0f;    // アタック時は約516Hz付近から172.3Hzへ急降下させる
    float bodySweepTime= 0.025f;  // 一瞬でピッチを落とす
    float bodyDecay    = 0.07f;   
    float bodyVolRatio = 0.63f * sinSet;    // 172.3Hzは強いエネルギーを持っているので、音量比率を高めに
    
    // --- 2. 倍音（カンカンしたヘッドの芯） ---
    float overtoneFreq  = 323.0f * random(0.98f, 1.02f); // ★データ2番目の「323.0 Hz (強さ: 0.55)」をそのまま指定！
    float overtoneDecay = 0.048f;  // 胴鳴り（172.3Hz）よりも早く減衰させて歯切れを良くする
    float overtoneVolRatio = 0.55f * sinSet;
    
    // --- 3. スナッピー（ホワイトノイズ） ---
    // データを見ると 1184.3 Hz 付近から徐々に強さが上がり、2.5kHz〜8kHzがエネルギーの塊になっています。
    //float noiseFilterCutoff = 2000.0f; // 2kHzあたりから上のノイズをまるごと通す
    float noiseDecay        = 0.17f;   // スナッピーの余韻
    float noiseVolRatio     = 3.5f;
    float noiseHPCutoff     = 2500.0f * random(0.9f, 1.1f);
    float noiseLPCutoff1    = 2800.0f * random(0.9f, 1.1f);
    float noiseLPCutoff2    = 4000.0f * random(0.9f, 1.1f);
    float noiseLPCutoff3    = 7000.0f * random(0.9f, 1.1f);

    // ==========================================
    // 各モジュールの構築とパッチング
    // ==========================================
    
    // --- 1. 胴鳴りパートの構築 ---
    bodyOsc = new Oscil(bodyFreq, baseAmp * bodyVolRatio, Waves.SINE);
    // 叩いた瞬間に音程を急激に下げる（アタック感の演出）
    bodyPitchEnv = new Line(bodySweepTime, bodyFreq * bodyPitchMod, bodyFreq);
    bodyPitchEnv.patch(bodyOsc.frequency);
    // 音量エンベロープ (Attack, Decay, Sustain, Release)
    bodyAmpEnv = new ADSR(1.0f, 0.001f, bodyDecay, 0.0f, 0.05f);
    
    bodyOsc.patch(bodyAmpEnv).patch(mixer);
    
    

    // --- 2. 倍音パートの構築 ---
    overtoneOsc = new Oscil(overtoneFreq, baseAmp * overtoneVolRatio, Waves.SINE);
    overtoneAmpEnv = new ADSR(1.0f, 0.002f, overtoneDecay, 0.0f, 0.05f);
    
    overtoneOsc.patch(overtoneAmpEnv).patch(mixer);
    
    

  // --- 3. スナッピー（Noise）パートの構築 ---
    snappyNoise = new Noise(baseAmp * noiseVolRatio, Noise.Tint.WHITE);
    
    // 4,500Hzに向かってなだらかに立ち上げるためのハイパス
    HighPassSP snappyHP = new HighPassSP(noiseHPCutoff, out.sampleRate());
    
    // 5,500Hzから20,000Hzに向けてなだらかに落とすためのローパス
    LowPassSP  snappyLP1 = new LowPassSP(noiseLPCutoff1, out.sampleRate());
    LowPassSP  snappyLP2 = new LowPassSP(noiseLPCutoff2, out.sampleRate());
    LowPassSP  snappyLP3 = new LowPassSP(noiseLPCutoff3, out.sampleRate());
    
    snappyAmpEnv = new ADSR(1.0f, 0.003f, noiseDecay, 0.0f, 0.05f);
    
    // ノイズを2つのフィルターに順番に通して（挟み撃ちにして）エンベロープへ送る
    snappyNoise.patch(snappyHP).patch(snappyLP1).patch(snappyLP2).patch(snappyLP3).patch(snappyAmpEnv).patch(mixer);
    
    

    // --- 4. マスターエンベロープ ---
    masterEnv = new ADSR(1.0f, 0.0001f, 0.0f, 1.0f, 0.01f);
    mixer.patch(masterEnv);
  }

  void noteOn(float duration) {
    masterEnv.noteOn();
    bodyAmpEnv.noteOn();
    overtoneAmpEnv.noteOn();
    snappyAmpEnv.noteOn();
    
    masterEnv.patch(shortOut);
  }

  void noteOff() {
    masterEnv.noteOff();
    bodyAmpEnv.noteOff();
    overtoneAmpEnv.noteOff();
    snappyAmpEnv.noteOff();
    
    masterEnv.unpatchAfterRelease(shortOut);
  }
}



// ==========================================
// 3. 新・ハイハット（Hi-Hat）のクラス (シンセサイズ方式)
// ==========================================
class HiHatInstrument implements Instrument {
  ADSR masterEnv;
  
  // 1. 金属共鳴（Metallic Core）パート用
  // グラフから厳選した主要な4つのピークを個別に発振・制御します
  //Oscil toneOsc1, toneOsc2, toneOsc3, toneOsc4;
  //ADSR  toneEnv1, toneEnv2, toneEnv3, toneEnv4;
  
  Oscil[] toneOscList;
  ADSR[] toneEnvList;
  
  // 2. 打撃摩擦（Noise）パート用
  Noise      hatNoise;
  //HighPassSP hatHP;
  ADSR       noiseEnv;

  HiHatInstrument(float masterVol, float velocity) {
    Summer mixer = new Summer();
    
    // 全体の音量調整（歪まないようにスネア等のバランスに合わせて調整してください）
    float baseAmp = masterVol * pow(constrain(velocity, 0, 127) / 127.0f, 0.9f) * 0.04f; 
    
    // ==========================================
    // パラメータ設定（ここをいじって音色をチューニングします！）
    // ==========================================
    
    // --- 1. 金属共鳴（チキチキ・カンカンした芯の成分） ---
    // FFT解析データの上位トップ4の強いピークをそのままスタンドアロンで指定
    /*
    float toneFreq1     = 7149.0f; // ★最強のセンターピーク（強さ: 1.00）
    float toneVolRatio1 = 1.00f;
    float toneDecay1    = 0.109f;  // 金属的なアタックのみを出すため、非常に短く
    
    float toneFreq2     = 7644.3f; // ★2番目に強い高域ピーク（強さ: 0.62）
    float toneVolRatio2 = 0.62f;
    float toneDecay2    = 0.040f;
    
    float toneFreq3     = 5964.7f; // ★3番目に強い中高域ピーク（強さ: 0.61）
    float toneVolRatio3 = 0.61f;
    float toneDecay3    = 0.060f;
    
    float toneFreq4     = 13006.1f;// ★超高域の金属的な鋭さを補うピーク（強さ: 0.52）
    float toneVolRatio4 = 0.52f;
    float toneDecay4    = 0.015f;  // 超高域は一瞬で消え去るように
    */
    
    float[][] toneList     = {
                               {21.53f, 0.0619f, 0.131f},
                               {7149.0f, 1.00f, 0.109f},
                               {7644.3f, 0.62f, 0.040f},
                               {5964.7f, 0.61f, 0.060f},
                               {13006.1f, 0.52f, 0.015f},
                               {5555.57f, 0.4585f, 0.118f},
                               {14610.28f, 0.1893f, 0.091f},
                               {4931.1f, 0.1834f, 0.109f}
                             };
    
    // --- 2. 打撃摩擦（シャリシャリ・シュワシュワした高域ノイズ） ---
    float noiseVolRatio = 6.00f * random(0.99f, 1.01f);   // ノイズのブレンド比率。シャリシャリ感を強めるなら上げる
    float noiseHPCutoff1 = 7000.0f * random(0.9f, 1.1f); // ★ハイパスのカットオフ（Hz）。
                                   // 8k〜10kHz以上にすると、ゴソゴソした中音域が消え、綺麗な「チッ」になります
    float noiseHPCutoff2 = 4000.0f * random(0.9f, 1.1f);
    float noiseLPCutoff1 = 5000.0f * random(0.9f, 1.1f);
    float noiseLPCutoff2 = 30000.0f * random(0.9f, 1.1f);
    //float noiseLPCutoff3 = 3000.0f;
    float noiseDecay    = 0.080f;  // クローズドハイハットの「余韻の長さ（キレ）」を左右する最重要パラメータ
    
    // ==========================================
    // 各モジュールの構築とパッチング
    // ==========================================
    
    // --- 1. 金属共鳴パートの構築 ---
    // 各正弦波の位相（Phase）をランダムにすることで、発音ごとの微妙な「なじみ」を生み出します
    /*
    toneOsc1 = new Oscil(toneFreq1, baseAmp * toneVolRatio1, Waves.SINE);
    toneOsc1.setPhase(random(1.0f));
    toneEnv1 = new ADSR(1.0f, 0.001f, toneDecay1, 0.0f, 0.01f);
    toneOsc1.patch(toneEnv1).patch(mixer);
    
    toneOsc2 = new Oscil(toneFreq2, baseAmp * toneVolRatio2, Waves.SINE);
    toneOsc2.setPhase(random(1.0f));
    toneEnv2 = new ADSR(1.0f, 0.001f, toneDecay2, 0.0f, 0.01f);
    toneOsc2.patch(toneEnv2).patch(mixer);
    
    toneOsc3 = new Oscil(toneFreq3, baseAmp * toneVolRatio3, Waves.SINE);
    toneOsc3.setPhase(random(1.0f));
    toneEnv3 = new ADSR(1.0f, 0.001f, toneDecay3, 0.0f, 0.01f);
    toneOsc3.patch(toneEnv3).patch(mixer);
    
    toneOsc4 = new Oscil(toneFreq4, baseAmp * toneVolRatio4, Waves.SINE);
    toneOsc4.setPhase(random(1.0f));
    toneEnv4 = new ADSR(1.0f, 0.001f, toneDecay4, 0.0f, 0.01f);
    toneOsc4.patch(toneEnv4).patch(mixer);
    */
    
    toneOscList = new Oscil[toneList.length];
    toneEnvList = new ADSR[toneList.length];
    
    for (int i = 0; i < toneList.length; i++){
      toneOscList[i] = new Oscil(toneList[i][0] * random(0.98f, 1.02f), baseAmp * toneList[i][1], Waves.SINE);
      toneOscList[i].setPhase(random(1.0f));
      toneEnvList[i] = new ADSR(1.0f, 0.001f, toneList[i][2]*0.5, 0.0f, 0.01f);
      toneOscList[i].patch(toneEnvList[i]).patch(mixer);
    }
    
    // --- 2. ノイズパートの構築 ---
    hatNoise = new Noise(baseAmp * noiseVolRatio, Noise.Tint.WHITE);
    HighPassSP hatHP1    = new HighPassSP(noiseHPCutoff1, out.sampleRate());
    HighPassSP hatHP2    = new HighPassSP(noiseHPCutoff2, out.sampleRate());
    LowPassSP hatLP1    = new LowPassSP(noiseLPCutoff1, out.sampleRate());
    LowPassSP hatLP2    = new LowPassSP(noiseLPCutoff2, out.sampleRate());
    //LowPassSP hatLP3    = new LowPassSP(noiseLPCutoff3, out.sampleRate());
    noiseEnv = new ADSR(1.0f, 0.001f, noiseDecay, 0.0f, 0.02f);
    
    // ホワイトノイズ -> ハイパスフィルター -> エンベロープ -> ミキサー
    hatNoise.patch(hatHP1).patch(hatHP2).patch(hatLP1).patch(hatLP2).patch(noiseEnv).patch(mixer);
    
    // --- 3. 全体の結合とマスターエンベロープ ---
    masterEnv = new ADSR(1.0f, 0.0001f, 0.0f, 1.0f, 0.02f); // リリースも極短に
    mixer.patch(masterEnv);
  }

  void noteOn(float duration) {
    masterEnv.noteOn();
    //toneEnv1.noteOn();
    //toneEnv2.noteOn();
    //toneEnv3.noteOn();
    //toneEnv4.noteOn();
    
    for (int i = 0; i<toneEnvList.length; i++ ) toneEnvList[i].noteOn();
    
    noiseEnv.noteOn();
    
    masterEnv.patch(shortOut);
  }

  void noteOff() {
    masterEnv.noteOff();
    //toneEnv1.noteOff();
    //toneEnv2.noteOff();
    //toneEnv3.noteOff();
    //toneEnv4.noteOff();
    
    for (int i = 0; i < toneEnvList.length; i++) {if (toneEnvList[i] != null) toneEnvList[i].noteOff();}
    
    
    noiseEnv.noteOff();
    
    masterEnv.unpatchAfterRelease(shortOut);
  }
}





//class SnareRoll implements Thread {
//  SnareRoll () {
    
//  }
//  void run() {
    
//  }
//}















/*
// ==========================================
// シンバル（Cymbal）のクラス
// ==========================================
class CymbalInstrument implements Instrument {
  ADSR masterEnv;
  ADSR[] harmonicEnvs;
  ADSR noiseEnv; // シャーンという余韻ノイズ用

  CymbalInstrument(float masterVol, float velocity) {
    Summer mixer = new Summer();
    
    // JSONデータの取得（ファイル名やキーは環境に合わせてください）
    float[][] cymbalLUT = utils.GetLUT("Cymbal"); 
    
    // 全体のマスターエンベロープ（余韻を長く残すためリリースは3.5秒）
    masterEnv = new ADSR(1.0f, 0.005f, 0.0f, 1.0f, 3.5f);
    mixer.patch(masterEnv);
    
    int numComponents = cymbalLUT.length;
    Oscil[] tones = new Oscil[numComponents];
    harmonicEnvs = new ADSR[numComponents];
    
    // 倍音成分が多く音が割れやすいため、ベース音量を少し控えめに設定
    float baseAmp = masterVol * pow(constrain(velocity, 0, 127) / 127.0f, 0.9f) * 1.0f;

    // ------------------------------------------------
    // 1. 金属のコア層（加算合成 ＋ FM/AM変調）
    // ------------------------------------------------
    for (int i = 0; i < numComponents; i++) {
      if (cymbalLUT[i] == null || cymbalLUT[i].length < 3) continue;
      
      float freq = cymbalLUT[i][0];
      float mag = cymbalLUT[i][1];
      float decayTime = cymbalLUT[i][2]; 
      
      // 高音域ほど音量が大きくなりすぎないよう微調整
      float amp = (baseAmp * mag) / numComponents/3;
      
      tones[i] = new Oscil(freq, amp, Waves.SINE);
      
      // アタック時の位相が揃って「バチッ」と鳴るのを防ぐ
      tones[i].setPhase(random(1.0f)); 
      
      // --- ① 非整数倍音の干渉をシミュレートするFM変調（ピッチ揺れ） ---
      // フルートよりも少し速く、ランダムな周期で揺らすことで金属の歪みを生む
      float fmSpeed = random(3.0f, 15.0f); 
      float vibDepthHz = freq * random(0.002f, 0.008f); 
      Oscil freqController = new Oscil(fmSpeed, vibDepthHz, Waves.SINE);
      freqController.offset.setLastValue(freq); 
      freqController.patch(tones[i].frequency);
      
      // --- ② 複雑なビートを生み出すAM変調（音量揺らぎ） ---
      // 倍音ごとに異なる周期で音量をウネウネさせる
      float amSpeed = random(1.0f, 8.0f);
      Oscil ampLFO = new Oscil(amSpeed, amp * random(0.3f, 0.7f), Waves.SINE);
      ampLFO.offset.setLastValue(amp);
      ampLFO.patch(tones[i].amplitude);
      
      // --- 個別エンベロープ（LUTのdecayTimeを適用） ---
      // decayTimeをそのままReleaseにも適用し、自然に消えさせる
      harmonicEnvs[i] = new ADSR(1.0f, 0.001f, decayTime/5, 0.0f, decayTime);
      
      tones[i].patch(harmonicEnvs[i]);
      harmonicEnvs[i].patch(mixer);
    }

    // ------------------------------------------------
    // 2. 超高域の余韻層（ノイズ ＋ ハイパスフィルター）
    // ------------------------------------------------
    // 「バーン」の後の「シャーン」という細かく散る粒子感をホワイトノイズで補完
    Noise cymbalNoise = new Noise(baseAmp * 0.10f, Noise.Tint.PINK);
    
    // 5000Hz以下の重たい成分を削り、チリチリとした超高域だけを残す
    HighPassSP hpf = new HighPassSP(5000f, out.sampleRate());
    
    // ノイズ層はゆっくり立ち上がり、コア層よりも長く残るように設定
    noiseEnv = new ADSR(1.0f, 0.02f, 1.5f, 0.3f, 3.5f); 
    
    cymbalNoise.patch(hpf).patch(noiseEnv).patch(mixer);
  }

  void noteOn(float duration) {
    masterEnv.noteOn();
    noiseEnv.noteOn();
    for (int i = 0; i < harmonicEnvs.length; i++) {
      if (harmonicEnvs[i] != null) harmonicEnvs[i].noteOn();
    }
    // 出力先（shortOut等）は環境に合わせて変更してください
    masterEnv.patch(out); 
  }

  void noteOff() {
    masterEnv.noteOff();
    noiseEnv.noteOff();
    for (int i = 0; i < harmonicEnvs.length; i++) {
      if (harmonicEnvs[i] != null) harmonicEnvs[i].noteOff();
    }
    masterEnv.unpatchAfterRelease(out);
  }
}*/
\end{lstlisting}

\subsection{Horn.pde}
\begin{lstlisting}[caption={Horn.pde}, label={lst:proc_horn}]
class HornInstrument implements Instrument {
  byte numHarmonics/* = 15*/;
  
  ADSR masterEnv;
  ADSR[] harmonicEnvs/* = new ADSR[numHarmonics]*/;
  
  Line[] wobbleLines;    // 6/12【追加】倍音ウナリ（AM）用のフェードライン
  float[] wobbleFactors; // 6/12【追加】倍音ごとの固有ウナリ比率の保持用
  
  // ==========================================
  // 7/5【追加1】ピッチエンベロープ用の変数
  // ==========================================
  Line[] pitchEnvLines;  // アタック時のピッチ急降下用
  float[] targetFreqs;   // 各倍音の本来の周波数（noteOnで計算に使うため保持）
  
  Line[] ampLines/* = new Line[numHarmonics]*/;
  float[] startAmps/* = new float[numHarmonics]*/;
  float[] endAmps/* = new float[numHarmonics]*/;
  
  ADSR breathEnv;
  
  
  // ==========================================
  // 【新規追加】 持続する空気の渦と管の共鳴（サブハーモニクス層）
  // ==========================================
  //ADSR sustainEnv;
  
  // 【追加】持続ノイズの音量変化（クレッシェンド/デクレッシェンド）用
  //Line sustainAmpLine;
  //Line sustainWobbleLine; // 6/12【追加】持続ノイズのウナリ用のフェードライン
  //float startSustainVol;
  //float endSustainVol;
  
  
  
  HornInstrument(float midiNote, float duration, float startVel, float endVel, float masterVol) {
    Summer mixer = new Summer();
    
    //float masterAttackTime = 0.2f - constrain(startVel, 0, 127) / 127.0f * 0.19f;
    float masterAttackTime = 0.1f;
    masterAttackTime = min(masterAttackTime, duration - 0.1f);
    if (masterAttackTime < 0.035f) masterAttackTime = 0.035f;
    
    // 【修正】masterEnv の一番最後の引数（リリース時間）を 0.3f から 0.4f に延ばす。
    // 内部の音が 0.3秒 かけて完全に消え去るのを待ってから、安全にアンパッチさせるための余裕（バッファ）です。
    // 【修正】masterEnv のリリース時間を、希望のフェードアウト時間（例：0.35秒）に設定
    masterEnv = new ADSR(1.0f, masterAttackTime, 0.0f, 1.0f, 0.13f);
    
    mixer.patch(masterEnv);
    
    
    // 変更6/12
    // 1. 倍音レシピを安全に取り出すための基準ベロシティ（0除算・ロード不全の防止）
    float profileVel = (startVel > 0) ? startVel : endVel;
    if (profileVel <= 0) profileVel = 64.0f; // どちらも0の場合のセーフティ
    
    // 2. 開始・終了のゲインファクターをフルートの特性（0.7乗カーブ）に基づいて独立計算
    float startGainFactor = pow(startVel / profileVel, 0.9f);
    float endGainFactor   = pow(endVel / profileVel, 0.9f);
    
    // 安全にロードされた基準プロファイルを取得
    float[][] startProfile = utils.GetDynamicProfile(midiNote, profileVel, "Horn", 34, 1.0f);
    
    float volFactor = masterVol * 0.017f;
    
    numHarmonics = (byte)startProfile.length;
    
    harmonicEnvs  = new ADSR[numHarmonics];
    ampLines      = new Line[numHarmonics];
    wobbleLines   = new Line[numHarmonics];    // 初期化を追加
    wobbleFactors = new float[numHarmonics];   // 初期化を追加
    startAmps     = new float[numHarmonics];
    endAmps       = new float[numHarmonics];
    
    // ==========================================
    // 7/5【追加2】配列の初期化
    // ==========================================
    pitchEnvLines = new Line[numHarmonics];
    targetFreqs   = new float[numHarmonics];
    
    
    
    //float fundamentalFreq = 440.0 * pow(2.0, (actualMidiNote - 69) / 12.0);
    byte baseWaveNum = 0;
    for (byte i = 1; i < numHarmonics; i++) {
      if (startProfile[i-1][0] < startProfile[i][0]) baseWaveNum = i;
    }
    
    
    // ==========================================
    // 【新規】倍音ごとの自然な振幅揺らぎ（変動率ベース）
    // ==========================================
    // Pythonで解析した変動率(%)を参考に、基音の振幅に対して何倍の揺れ(±)を許容するかを設定。
    // 実測データの変動率(%)から数式 [ Variation / 500 ] によって導出された、
    // 倍音ごとの自然な振幅揺らぎ（AM変調）の深度設定。
    //float[] wobbleDepths = {0.1250f, 0.1616f, 0.2078f, 0.2340f, 0.2114f, 0.2742f, 0.3716f, 0.2946f, 0.7420f, 0.8796f, 
    //                        0.7850f, 0.5070f, 0.6242f, 0.7654f, 0.6702f, 0.3898f, 0.8850f, 0.5942f, 0.3152f, 0.5030f};
    
    // ==========================================
    // 【新規】この音符全体の「呼吸のスピード」を決定
    // forループの外で1つだけ定義し、全倍音で共有することで相殺を防ぐ
    // ==========================================
    float globalBreathSpeed = random(2.5f, 3.5f); // 生楽器の自然なビブラート周期
    
    
    


    // 1. 加算合成パート（7つのサイン波）
    for (int i = 0; i < numHarmonics; i++) {
      /*
      int h = i + 1;
      // あなたがPythonで導き出した厳密なインハーモニシティ公式を適用
      float B = 0.0005f; // インハーモニシティ係数
      float inharmonicity = sqrt(1.0f + B * h * h);
      float targetFreq = f0 * h * inharmonicity;
      */
      float targetFreq = startProfile[i][0];
      
      Oscil osc = new Oscil(targetFreq, 0.0f, Waves.SINE);
      
      // 【追加】初期位相を 0.0 〜 1.0 (0度〜360度) の間で完全にランダムに散らす
      // これにより「カチッ」とした機械的なアタック音が消え、音割れ（ピーク）も防げる
      osc.setPhase(random(1.0f));
      
      
      //startAmps[i] = startProfile[i][1] * volFactor;
      //endAmps[i] = endProfile[i] * volFactor;
        // 聴覚特性（0.7乗）に完全同期させた、開始時から終了時への音量変化比率を算出
        //float velRatio = pow(endVel / (startVel + 0.0001f), 0.7f);
      // 【確定修正】すでにvolFactorが掛けられたstartAmpsに対して、正しいべき乗比率を乗算する
      //endAmps[i]   = startAmps[i] * pow(endVel / (startVel + 0.0001f), 0.7f);
      
      // 6/12追加
      // 独立した開始・終了音量を厳密に算出
      float baseAmp = startProfile[i][1] * volFactor;
      startAmps[i] = baseAmp * startGainFactor;
      endAmps[i]   = baseAmp * endGainFactor;
      
      
      
        // 通常のビブラート（サイン波）
        
        // ==========================================
        // ① 極小のピッチ揺れ (FM変調) の復活
        // 以前の0.0132fを「0.003f」に激減させ、音痴にならない程度の「空気の揺らぎ」を作る
        // ==========================================
        float vibDepthHz = targetFreq * 0.001f; 
        Oscil freqController = new Oscil(globalBreathSpeed, vibDepthHz, Waves.SINE);
        freqController.offset.setLastValue(targetFreq); 
        
        
        //freqController.patch(osc.frequency);
        
        // ==========================================
        // 7/5【追加3】本来の周波数を保存し、Pitch用のLineを接続
        // ==========================================
        targetFreqs[i] = targetFreq; // noteOnで使うために記憶しておく
        
        //pitchEnvLines[i] = new Line();
        //pitchEnvLines[i].patch(osc.frequency); // オシレーターの周波数に加算接続
        // 【変更点】7/5
        // 以前あった freqController.offset.setLastValue(targetFreq); を削除します。
        // 代わりに、Pitch用のLineをLFOの「offset（基準周波数）」に直接パッチします！
        pitchEnvLines[i] = new Line();
        pitchEnvLines[i].patch(freqController.offset); 
        
        // LFO(基準値がLineで動く)を、メインオシレーターの周波数にパッチ
        freqController.patch(osc.frequency);
        
        // ==========================================
        // ② 呼吸に同期した音量揺らぎ (AM変調)
        // ==========================================
        // この倍音の基準音量に対して、配列で設定した割合だけ揺らす
        //float wobbleAmp = startAmps[i] * wobbleDepths[i]; 
        //float wobbleAmp = startAmps[i] * random(0.1f, 0.3f);
        
        // 1.5Hz 〜 3.5Hz の間で、倍音ごとにランダムなスピードで揺らす
        // 【修正】ランダムではなく、統一された globalBreathSpeed を使う！
        //Oscil ampLFO = new Oscil(globalBreathSpeed, wobbleAmp, Waves.SINE);
        
        // 位相（スタート位置）は少しだけバラけさせると、音が立体的になります
        //ampLFO.setPhase(random(1.0f)); 
        
        // ampLFOを osc.amplitude にパッチ(追加)する。
        // Minimでは複数のUGenを同じ入力にパッチすると「加算(Sum)」されるため、
        // ampLines の基準値に対して、ampLFO が ±wobbleAmp の揺れを足し引きしてくれます。
        //ampLFO.patch(osc.amplitude);
        
        // --- 音量揺らぎ（AM変調）のフェード構造化 ---
        wobbleFactors[i] = random(0.12f, 0.17f); 
        //wobbleFactors[i] = wobbleDepths[i]; 
        wobbleLines[i] = new Line();
        Oscil ampLFO = new Oscil(globalBreathSpeed, 0.0f, Waves.SINE); // 初期振幅は0
        ampLFO.setPhase(random(1.0f)); 
        
        wobbleLines[i].patch(ampLFO.amplitude); // 新設LineをLFOの振幅へ接続
        ampLFO.patch(osc.amplitude);
      
      
      ampLines[i] = new Line();
      ampLines[i].patch(osc.amplitude); 
      
      float attackTime = (i == baseWaveNum) ? 0.058f : 0.395f; 
      
      // 【修正】内部のサイン波のリリースを 0.3f から 10.0f（10秒）にして、急激な角を作らせない
      harmonicEnvs[i] = new ADSR(1.0f, attackTime/3, 1.0f, 0.7f, 10.0f);
      osc.patch(harmonicEnvs[i]);
      harmonicEnvs[i].patch(mixer);
    }
    
    
    // ==========================================
    // 2. 息のノイズ（Chiff）パートの実測値アップデート
    // =========================================
    // 0.0〜1.0 のベロシティ割合を計算（ノイズの音量調整用）
    float startNorm = constrain(startVel, 0, 127) / 127.0f;
    
    // 中低音の破裂音（ドスッという音）になるため、少しだけ音量を上げます（0.01f -> 0.02f）
    float attackNoiseVol = masterVol * 0.003f * pow(startNorm, 0.9f); 
    Noise breathNoise = new Noise(attackNoiseVol, Noise.Tint.WHITE);
    
    // 【修正】ハイパス(HP)からバンドパス(BP)に変更。
    // 3500Hz付近の「シュッ」という音だけを残し、耳障りな高音を削る。
    // 【修正】実測データに基づき、3500Hzから 581.4Hz に変更！
    // 473Hzの山も一緒に拾うため、レゾナンスを 0.3 から 0.15 に下げて帯域を広げる
    MoogFilter breathFilter = new MoogFilter(381.4f, 0.15f);
    breathFilter.type = MoogFilter.Type.BP;
    
    breathEnv = new ADSR(1.0f, 0.01f, 0.08f, 0.0f, 0.0f);
    
    breathNoise.patch(breathFilter);
    breathFilter.patch(breathEnv);
    breathEnv.patch(mixer);
    
    
    /*
    // ==========================================
    // 3. 【新規追加】持続する空気の渦と管の共鳴（サブハーモニクス層）
    // ==========================================
    float endNorm = constrain(endVel, 0, 127) / 127.0f; // 【追加】終了時のベロシティ割合
    
    // 【修正】開始時と終了時のノイズ音量を計算して保持する
    startSustainVol = masterVol * 0.005f * pow(startNorm, 0.9f);
    endSustainVol   = masterVol * 0.005f * pow(endNorm, 0.9f);
    
    // ピンクノイズで低音の「フオォォ」という空気感を出す
    // 【修正】Noiseの初期音量は0.0fにし、音量制御はLineに任せる
    Noise sustainNoise = new Noise(0.0f, Noise.Tint.PINK);
    
    // 【追加】持続ノイズのボリュームつまみにLineを接続
    sustainAmpLine = new Line();
    sustainAmpLine.patch(sustainNoise.amplitude);
    
    // ==========================================
    // 【新規】持続ノイズ自体にも「息の乱れ」を適用する
    // サイン波の第1倍音と同等のスピードと深さでノイズを揺らす
    // ==========================================
    //float noiseWobbleAmp = startSustainVol * 0.20f; // 20%揺らす
    //Oscil sustainWobble = new Oscil(random(1.5f, 3.0f), noiseWobbleAmp, Waves.SINE);
    //sustainWobble.setPhase(random(1.0f));
    //sustainWobble.patch(sustainNoise.amplitude); // Lineの音量にLFOの揺れを加算
    // ==========================================
    
    // 6/12変更
    // --- 持続ノイズの息の乱れも完全にフェードさせる ---
    sustainWobbleLine = new Line();
    Oscil sustainWobble = new Oscil(random(1.5f, 3.0f), 0.0f, Waves.SINE); // 初期振幅は0
    sustainWobble.setPhase(random(1.0f));
    
    sustainWobbleLine.patch(sustainWobble.amplitude); // 新設LineをノイズLFOの振幅へ接続
    sustainWobble.patch(sustainNoise.amplitude);      // LFO出力をノイズ振幅へ（加算）
    
    // ローパスフィルターで高音を削りつつ、2500Hz付近に共鳴(レゾナンス0.5)の山を作る
    MoogFilter sustainFilter = new MoogFilter(2500f, 0.5f);
    sustainFilter.type = MoogFilter.Type.LP;
    
    // 音が鳴っている間ずっと持続するエンベロープ（サイン波のアタックに合わせる）
    // 【修正】持続ノイズのリリースも 0.3f から 10.0f（10秒）にする
    sustainEnv = new ADSR(1.0f, 0.058f, 0.0f, 1.0f, 10.0f);
    
    sustainNoise.patch(sustainFilter);
    sustainFilter.patch(sustainEnv);
    sustainEnv.patch(mixer);
    */
  }
  
  
  /*
  void noteOn(float dur) {
    masterEnv.noteOn();
    
    // ====================================================
    // 【修正】Lineの寿命に、リリース時間（0.4秒）分の猶予を足す！
    // これにより、減衰中もLineが生き残り、オシレーターの音量が突然0.0fにリセットされるのを防ぐ
    // ====================================================
    // 【微調整】masterEnvのリリース時間(0.35f)より少し長ければOK
    float safeDur = dur + 0.5f;
    
    for(int i=0; i<numHarmonics; i++) {
      harmonicEnvs[i].noteOn();
      ampLines[i].activate(safeDur, startAmps[i], endAmps[i]);
    }
    
    breathEnv.noteOn();
    sustainEnv.noteOn(); // 【追加】持続ノイズをオン
    sustainAmpLine.activate(safeDur, startSustainVol, endSustainVol); // 【追加】発音と同時に、持続ノイズの音量変化をスタートさせる
    
    masterEnv.patch(out); 
  }*/
  // 6/12変更
  void noteOn(float dur) {
    masterEnv.noteOn();
    float safeDur = dur + 0.5f;
    
    // 7/5【追加4】トランペット/ホルン特有の「アタック時のピッチ降下」パラメータ
    float sweepTime = 0.12f; // 変化にかかる時間（0.04秒）
    float pitchMod  = 0.05f; // 高くなる割合（0.03 = 本来の周波数の3%増しからスタート）
    
    for(int i = 0; i < numHarmonics; i++) {
      harmonicEnvs[i].noteOn();
      
      // 1. 各倍音のベース音量をフェード
      ampLines[i].activate(safeDur, startAmps[i], endAmps[i]);
      
      // 2. 通常ビブラート時のみ、ウナリ（AM）の深さも比例フェードさせる
      if (wobbleLines[i] != null) {
        float factor = wobbleFactors[i];
        wobbleLines[i].activate(safeDur, startAmps[i] * factor, endAmps[i] * factor);
      }
      
      // 3. 7/5【新規追加】アタック時のピッチエンベロープ（周波数の急降下）
      // 「本来の周波数の3%分上乗せした値」から「0（上乗せなし）」へと一直線に下げる
      //if (pitchEnvLines[i] != null) {
      //  float baseFreq = targetFreqs[i];
      //  pitchEnvLines[i].activate(sweepTime, baseFreq * pitchMod, 0.0f);
      //}
      // 3. 7/5【新規追加】アタック時のピッチエンベロープ（周波数の急降下）
      if (pitchEnvLines[i] != null) {
        float baseFreq = targetFreqs[i];
        
        // スタート地点の周波数: 本来の周波数に pitchMod(%) を上乗せした値
        // 例: pitchMod が 0.03 なら 3%増し、2.0 なら 200%増し（3倍）
        float startFreq = baseFreq * (1.0f + pitchMod);
        
        // startFreq から baseFreq(本来の周波数) へ、一直線に下げる
        pitchEnvLines[i].activate(sweepTime, startFreq, baseFreq);
      }
    }
    
    breathEnv.noteOn();
    //sustainEnv.noteOn(); 
    
    // 3. 持続ノイズの全体音量と、そのウナリの深さを同時に追従フェード
    //sustainAmpLine.activate(safeDur, startSustainVol, endSustainVol); 
    //sustainWobbleLine.activate(safeDur, startSustainVol * 0.20f, endSustainVol * 0.20f); 
    
    masterEnv.patch(longOut); 
  }
  
  
  
  void noteOff() {
    masterEnv.noteOff();
    for(int i=0; i<numHarmonics; i++) harmonicEnvs[i].noteOff();
    
    // 【削除】breathEnv は既に無音なので noteOff を呼ばない（空撃ちグリッチ防止）
    // breathEnv.noteOff();
    
    //sustainEnv.noteOff(); // 【追加】持続ノイズをオフ
    masterEnv.unpatchAfterRelease(longOut); 
  }
}
\end{lstlisting}

\subsection{Trumpet.pde}
\begin{lstlisting}[caption={Trumpet.pde}, label={lst:proc_trumpet}]
class TrumpetInstrument implements Instrument {
  byte numHarmonics/* = 15*/;
  
  ADSR masterEnv;
  ADSR[] harmonicEnvs/* = new ADSR[numHarmonics]*/;
  
  Line[] wobbleLines;    // 6/12【追加】倍音ウナリ（AM）用のフェードライン
  float[] wobbleFactors; // 6/12【追加】倍音ごとの固有ウナリ比率の保持用
  
  // ==========================================
  // 7/5【追加1】ピッチエンベロープ用の変数
  // ==========================================
  Line[] pitchEnvLines;  // アタック時のピッチ急降下用
  float[] targetFreqs;   // 各倍音の本来の周波数（noteOnで計算に使うため保持）
  
  Line[] ampLines/* = new Line[numHarmonics]*/;
  float[] startAmps/* = new float[numHarmonics]*/;
  float[] endAmps/* = new float[numHarmonics]*/;
  
  ADSR breathEnv;
  
  
  // ==========================================
  // 【新規追加】 持続する空気の渦と管の共鳴（サブハーモニクス層）
  // ==========================================
  //ADSR sustainEnv;
  
  // 【追加】持続ノイズの音量変化（クレッシェンド/デクレッシェンド）用
  //Line sustainAmpLine;
  //Line sustainWobbleLine; // 6/12【追加】持続ノイズのウナリ用のフェードライン
  //float startSustainVol;
  //float endSustainVol;
  
  
  
  TrumpetInstrument(float midiNote, float duration, float startVel, float endVel, float masterVol) {
    Summer mixer = new Summer();
    
    float masterAttackTime = 0.1f - constrain(startVel, 0, 127) / 127.0f * 0.8f;
    masterAttackTime = min(masterAttackTime, duration - 0.05f);
    if (masterAttackTime < 0.035f) masterAttackTime = 0.035f;
    
    // 【修正】masterEnv の一番最後の引数（リリース時間）を 0.3f から 0.4f に延ばす。
    // 内部の音が 0.3秒 かけて完全に消え去るのを待ってから、安全にアンパッチさせるための余裕（バッファ）です。
    // 【修正】masterEnv のリリース時間を、希望のフェードアウト時間（例：0.35秒）に設定
    masterEnv = new ADSR(1.0f, masterAttackTime, 0.0f, 1.0f, 0.17f);
    
    mixer.patch(masterEnv);
    
    
    // 変更6/12
    // 1. 倍音レシピを安全に取り出すための基準ベロシティ（0除算・ロード不全の防止）
    float profileVel = (startVel > 0) ? startVel : endVel;
    if (profileVel <= 0) profileVel = 64.0f; // どちらも0の場合のセーフティ
    
    // 2. 開始・終了のゲインファクターをフルートの特性（0.7乗カーブ）に基づいて独立計算
    float startGainFactor = pow(startVel / profileVel, 0.9f);
    float endGainFactor   = pow(endVel / profileVel, 0.9f);
    
    // 安全にロードされた基準プロファイルを取得
    float[][] startProfile = utils.GetDynamicProfile(midiNote, profileVel, "Trumpet", 52, 1.0f);
    
    float volFactor = masterVol * 0.018f;
    
    numHarmonics = (byte)startProfile.length;
    
    harmonicEnvs  = new ADSR[numHarmonics];
    ampLines      = new Line[numHarmonics];
    wobbleLines   = new Line[numHarmonics];    // 初期化を追加
    wobbleFactors = new float[numHarmonics];   // 初期化を追加
    startAmps     = new float[numHarmonics];
    endAmps       = new float[numHarmonics];
    
    // ==========================================
    // 7/5【追加2】配列の初期化
    // ==========================================
    pitchEnvLines = new Line[numHarmonics];
    targetFreqs   = new float[numHarmonics];
    
    
    
    
    //float fundamentalFreq = 440.0 * pow(2.0, (actualMidiNote - 69) / 12.0);
    byte baseWaveNum = 0;
    for (byte i = 1; i < numHarmonics; i++) {
      if (startProfile[i-1][0] < startProfile[i][0]) baseWaveNum = i;
    }
    
    
    // ==========================================
    // 【新規】倍音ごとの自然な振幅揺らぎ（変動率ベース）
    // ==========================================
    // Pythonで解析した変動率(%)を参考に、基音の振幅に対して何倍の揺れ(±)を許容するかを設定。
    // 実測データの変動率(%)から数式 [ Variation / 500 ] によって導出された、
    // 倍音ごとの自然な振幅揺らぎ（AM変調）の深度設定。
    //float[] wobbleDepths = {0.1250f, 0.1616f, 0.2078f, 0.2340f, 0.2114f, 0.2742f, 0.3716f, 0.2946f, 0.7420f, 0.8796f, 
    //                        0.7850f, 0.5070f, 0.6242f, 0.7654f, 0.6702f, 0.3898f, 0.8850f, 0.5942f, 0.3152f, 0.5030f};
    
    // ==========================================
    // 【新規】この音符全体の「呼吸のスピード」を決定
    // forループの外で1つだけ定義し、全倍音で共有することで相殺を防ぐ
    // ==========================================
    float globalBreathSpeed = random(4.0f, 5.0f); // 生楽器の自然なビブラート周期
    
    
    


    // 1. 加算合成パート（7つのサイン波）
    for (int i = 0; i < numHarmonics; i++) {
      /*
      int h = i + 1;
      // あなたがPythonで導き出した厳密なインハーモニシティ公式を適用
      float B = 0.0005f; // インハーモニシティ係数
      float inharmonicity = sqrt(1.0f + B * h * h);
      float targetFreq = f0 * h * inharmonicity;
      */
      float targetFreq = startProfile[i][0];
      
      Oscil osc = new Oscil(targetFreq, 0.0f, Waves.SINE);
      
      // 【追加】初期位相を 0.0 〜 1.0 (0度〜360度) の間で完全にランダムに散らす
      // これにより「カチッ」とした機械的なアタック音が消え、音割れ（ピーク）も防げる
      osc.setPhase(random(1.0f));
      
      
      
      
      //startAmps[i] = startProfile[i][1] * volFactor;
      //endAmps[i] = endProfile[i] * volFactor;
        // 聴覚特性（0.7乗）に完全同期させた、開始時から終了時への音量変化比率を算出
        //float velRatio = pow(endVel / (startVel + 0.0001f), 0.7f);
      // 【確定修正】すでにvolFactorが掛けられたstartAmpsに対して、正しいべき乗比率を乗算する
      //endAmps[i]   = startAmps[i] * pow(endVel / (startVel + 0.0001f), 0.7f);
      
      // 6/12追加
      // 独立した開始・終了音量を厳密に算出
      float baseAmp = startProfile[i][1] * volFactor;
      startAmps[i] = baseAmp * startGainFactor;
      endAmps[i]   = baseAmp * endGainFactor;
      
      
      
        // 通常のビブラート（サイン波）
        
        // ==========================================
        // ① 極小のピッチ揺れ (FM変調) の復活
        // 以前の0.0132fを「0.003f」に激減させ、音痴にならない程度の「空気の揺らぎ」を作る
        // ==========================================
        float vibDepthHz = targetFreq * 0.0007f; 
        Oscil freqController = new Oscil(globalBreathSpeed, vibDepthHz, Waves.SINE);
        
        
        
        //freqController.offset.setLastValue(targetFreq); 
        
        // ==========================================
        // 7/5【追加3】本来の周波数を保存し、Pitch用のLineを接続
        // ==========================================
        targetFreqs[i] = targetFreq; // noteOnで使うために記憶しておく
        
        //pitchEnvLines[i] = new Line();
        //pitchEnvLines[i].patch(osc.frequency); // オシレーターの周波数に加算接続
        // 【変更点】7/5
        // 以前あった freqController.offset.setLastValue(targetFreq); を削除します。
        // 代わりに、Pitch用のLineをLFOの「offset（基準周波数）」に直接パッチします！
        pitchEnvLines[i] = new Line();
        pitchEnvLines[i].patch(freqController.offset); 
        
        // LFO(基準値がLineで動く)を、メインオシレーターの周波数にパッチ
        freqController.patch(osc.frequency);
        
        
        
        // ==========================================
        // ② 呼吸に同期した音量揺らぎ (AM変調)
        // ==========================================
        // この倍音の基準音量に対して、配列で設定した割合だけ揺らす
        //float wobbleAmp = startAmps[i] * wobbleDepths[i]; 
        //float wobbleAmp = startAmps[i] * random(0.1f, 0.3f);
        
        // 1.5Hz 〜 3.5Hz の間で、倍音ごとにランダムなスピードで揺らす
        // 【修正】ランダムではなく、統一された globalBreathSpeed を使う！
        //Oscil ampLFO = new Oscil(globalBreathSpeed, wobbleAmp, Waves.SINE);
        
        // 位相（スタート位置）は少しだけバラけさせると、音が立体的になります
        //ampLFO.setPhase(random(1.0f)); 
        
        // ampLFOを osc.amplitude にパッチ(追加)する。
        // Minimでは複数のUGenを同じ入力にパッチすると「加算(Sum)」されるため、
        // ampLines の基準値に対して、ampLFO が ±wobbleAmp の揺れを足し引きしてくれます。
        //ampLFO.patch(osc.amplitude);
        
        // --- 音量揺らぎ（AM変調）のフェード構造化 ---
        wobbleFactors[i] = random(0.10f, 0.20f); 
        //wobbleFactors[i] = wobbleDepths[i]; 
        wobbleLines[i] = new Line();
        Oscil ampLFO = new Oscil(globalBreathSpeed, 0.0f, Waves.SINE); // 初期振幅は0
        ampLFO.setPhase(random(1.0f)); 
        
        wobbleLines[i].patch(ampLFO.amplitude); // 新設LineをLFOの振幅へ接続
        ampLFO.patch(osc.amplitude);
      
      
      ampLines[i] = new Line();
      ampLines[i].patch(osc.amplitude); 
      
      float attackTime = (i == baseWaveNum) ? 0.058f : 0.395f; 
      
      // 【修正】内部のサイン波のリリースを 0.3f から 10.0f（10秒）にして、急激な角を作らせない
      harmonicEnvs[i] = new ADSR(1.0f, attackTime/3, 1.0f, 0.7f, 10.0f);
      osc.patch(harmonicEnvs[i]);
      harmonicEnvs[i].patch(mixer);
    }
    
    
    // ==========================================
    // 2. 息のノイズ（Chiff）パートの実測値アップデート
    // =========================================
    // 0.0〜1.0 のベロシティ割合を計算（ノイズの音量調整用）
    float startNorm = constrain(startVel, 0, 127) / 127.0f;
    
    // 中低音の破裂音（ドスッという音）になるため、少しだけ音量を上げます（0.01f -> 0.02f）
    float attackNoiseVol = masterVol * 0.01f * pow(startNorm, 0.9f); 
    Noise breathNoise = new Noise(attackNoiseVol, Noise.Tint.WHITE);
    
    // 【修正】ハイパス(HP)からバンドパス(BP)に変更。
    // 3500Hz付近の「シュッ」という音だけを残し、耳障りな高音を削る。
    // 【修正】実測データに基づき、3500Hzから 581.4Hz に変更！
    // 473Hzの山も一緒に拾うため、レゾナンスを 0.3 から 0.15 に下げて帯域を広げる
    MoogFilter breathFilter = new MoogFilter(581.4f, 0.15f);
    breathFilter.type = MoogFilter.Type.BP;
    
    breathEnv = new ADSR(1.0f, 0.01f, 0.08f, 0.0f, 0.0f);
    
    breathNoise.patch(breathFilter);
    breathFilter.patch(breathEnv);
    breathEnv.patch(mixer);
    
    
    /*
    // ==========================================
    // 3. 【新規追加】持続する空気の渦と管の共鳴（サブハーモニクス層）
    // ==========================================
    float endNorm = constrain(endVel, 0, 127) / 127.0f; // 【追加】終了時のベロシティ割合
    
    // 【修正】開始時と終了時のノイズ音量を計算して保持する
    startSustainVol = masterVol * 0.005f * pow(startNorm, 0.9f);
    endSustainVol   = masterVol * 0.005f * pow(endNorm, 0.9f);
    
    // ピンクノイズで低音の「フオォォ」という空気感を出す
    // 【修正】Noiseの初期音量は0.0fにし、音量制御はLineに任せる
    Noise sustainNoise = new Noise(0.0f, Noise.Tint.PINK);
    
    // 【追加】持続ノイズのボリュームつまみにLineを接続
    sustainAmpLine = new Line();
    sustainAmpLine.patch(sustainNoise.amplitude);
    
    // ==========================================
    // 【新規】持続ノイズ自体にも「息の乱れ」を適用する
    // サイン波の第1倍音と同等のスピードと深さでノイズを揺らす
    // ==========================================
    //float noiseWobbleAmp = startSustainVol * 0.20f; // 20%揺らす
    //Oscil sustainWobble = new Oscil(random(1.5f, 3.0f), noiseWobbleAmp, Waves.SINE);
    //sustainWobble.setPhase(random(1.0f));
    //sustainWobble.patch(sustainNoise.amplitude); // Lineの音量にLFOの揺れを加算
    // ==========================================
    
    // 6/12変更
    // --- 持続ノイズの息の乱れも完全にフェードさせる ---
    sustainWobbleLine = new Line();
    Oscil sustainWobble = new Oscil(random(1.5f, 3.0f), 0.0f, Waves.SINE); // 初期振幅は0
    sustainWobble.setPhase(random(1.0f));
    
    sustainWobbleLine.patch(sustainWobble.amplitude); // 新設LineをノイズLFOの振幅へ接続
    sustainWobble.patch(sustainNoise.amplitude);      // LFO出力をノイズ振幅へ（加算）
    
    // ローパスフィルターで高音を削りつつ、2500Hz付近に共鳴(レゾナンス0.5)の山を作る
    MoogFilter sustainFilter = new MoogFilter(2500f, 0.5f);
    sustainFilter.type = MoogFilter.Type.LP;
    
    // 音が鳴っている間ずっと持続するエンベロープ（サイン波のアタックに合わせる）
    // 【修正】持続ノイズのリリースも 0.3f から 10.0f（10秒）にする
    sustainEnv = new ADSR(1.0f, 0.058f, 0.0f, 1.0f, 10.0f);
    
    sustainNoise.patch(sustainFilter);
    sustainFilter.patch(sustainEnv);
    sustainEnv.patch(mixer);
    */
  }
  
  
  /*
  void noteOn(float dur) {
    masterEnv.noteOn();
    
    // ====================================================
    // 【修正】Lineの寿命に、リリース時間（0.4秒）分の猶予を足す！
    // これにより、減衰中もLineが生き残り、オシレーターの音量が突然0.0fにリセットされるのを防ぐ
    // ====================================================
    // 【微調整】masterEnvのリリース時間(0.35f)より少し長ければOK
    float safeDur = dur + 0.5f;
    
    for(int i=0; i<numHarmonics; i++) {
      harmonicEnvs[i].noteOn();
      ampLines[i].activate(safeDur, startAmps[i], endAmps[i]);
    }
    
    breathEnv.noteOn();
    sustainEnv.noteOn(); // 【追加】持続ノイズをオン
    sustainAmpLine.activate(safeDur, startSustainVol, endSustainVol); // 【追加】発音と同時に、持続ノイズの音量変化をスタートさせる
    
    masterEnv.patch(out); 
  }*/
  // 6/12変更
  void noteOn(float dur) {
    masterEnv.noteOn();
    float safeDur = dur + 0.5f;
    
    // 7/5【追加4】トランペット/ホルン特有の「アタック時のピッチ降下」パラメータ
    float sweepTime = 0.07f; // 変化にかかる時間（0.04秒）
    float pitchMod  = 0.03f; // 高くなる割合（0.03 = 本来の周波数の3%増しからスタート）
    
    for(int i = 0; i < numHarmonics; i++) {
      harmonicEnvs[i].noteOn();
      
      // 1. 各倍音のベース音量をフェード
      ampLines[i].activate(safeDur, startAmps[i], endAmps[i]);
      
      // 2. 通常ビブラート時のみ、ウナリ（AM）の深さも比例フェードさせる
      if (wobbleLines[i] != null) {
        float factor = wobbleFactors[i];
        wobbleLines[i].activate(safeDur, startAmps[i] * factor, endAmps[i] * factor);
      }
      
      // 3. 7/5【新規追加】アタック時のピッチエンベロープ（周波数の急降下）
      // 「本来の周波数の3%分上乗せした値」から「0（上乗せなし）」へと一直線に下げる
      //if (pitchEnvLines[i] != null) {
      //  float baseFreq = targetFreqs[i];
      //  pitchEnvLines[i].activate(sweepTime, baseFreq * pitchMod, 0.0f);
      //}
      // 3. 7/5【新規追加】アタック時のピッチエンベロープ（周波数の急降下）
      if (pitchEnvLines[i] != null) {
        float baseFreq = targetFreqs[i];
        
        // スタート地点の周波数: 本来の周波数に pitchMod(%) を上乗せした値
        // 例: pitchMod が 0.03 なら 3%増し、2.0 なら 200%増し（3倍）
        float startFreq = baseFreq * (1.0f + pitchMod);
        
        // startFreq から baseFreq(本来の周波数) へ、一直線に下げる
        pitchEnvLines[i].activate(sweepTime, startFreq, baseFreq);
      }
    }
    
    breathEnv.noteOn();
    //sustainEnv.noteOn(); 
    
    // 3. 持続ノイズの全体音量と、そのウナリの深さを同時に追従フェード
    //sustainAmpLine.activate(safeDur, startSustainVol, endSustainVol); 
    //sustainWobbleLine.activate(safeDur, startSustainVol * 0.20f, endSustainVol * 0.20f); 
    
    masterEnv.patch(longOut); 
  }
  
  
  
  void noteOff() {
    masterEnv.noteOff();
    for(int i=0; i<numHarmonics; i++) harmonicEnvs[i].noteOff();
    
    // 【削除】breathEnv は既に無音なので noteOff を呼ばない（空撃ちグリッチ防止）
    // breathEnv.noteOff();
    
    //sustainEnv.noteOff(); // 【追加】持続ノイズをオフ
    masterEnv.unpatchAfterRelease(longOut); 
  }
}
\end{lstlisting}
